\part{Holografía Modular Triádica}

\chapter{Toro aritmético nonádico de caminos y doble evaluación}
\label{chap:soporte-aritmetico-rev7}
\HMTClaim{HMT-R-APP-SINGLE-SUPPORT}
\HMTClaim{HMT-R-APP-DOUBLE-EVALUATION}
\HMTDemoteHeadings
\def\HMTAPPFormal{1}
\input{sections/hmt/03_app_ortograma}
\let\HMTAPPFormal\undefined
\def\HMTAPPDevelopments{1}
\input{sections/hmt/03_app_ortograma}
\let\HMTAPPDevelopments\undefined
\input{sections/hmt/03d_esqueleto_ciclotomico_app}
\HMTRestoreHeadings

\chapter{Coordenada ternaria orientada y regímenes cuadráticos}
\label{chap:trit-rev7}
\HMTClaim{HMT-R-TRIT-ORIENTATION}
\HMTClaim{HMT-R-TRIT-QUADRATIC}
\HMTDemoteHeadings
\def\HMTTRITFormal{1}
\input{sections/hmt/02_trit_geometrias}
\let\HMTTRITFormal\undefined
\def\HMTTRITDevelopments{1}
\input{sections/hmt/02_trit_geometrias}
\let\HMTTRITDevelopments\undefined
\HMTRestoreHeadings

\chapter{Núcleo topológico de fase: transporte categórico y memoria}
\label{chap:tpk-rev7}
\HMTClaim{HMT-R-TPK-ENRICHED}
\HMTClaim{HMT-R-TPK-MODULAR-READERS}
\HMTClaim{HMT-R-TPK-TEMPORAL-SELECTOR}
\HMTClaim{HMT-R-TPK-CARDINAL-TRANSPORT}
\HMTClaim{HMT-R-TPK-C12}
\HMTDemoteHeadings
\input{sections/hmt/03c_nucleo_transporte}
\input{sections/hmt/03e_funtor_fase_longitud}
\input{sections/hmt/03b_alfabeto_factorizacion}
\input{sections/hmt/20_arquitectura_operatoria_tpk_actualizada}
\input{sections/hmt/03g_ciclo_dodecafasico_tpk_rev7}
\HMTRestoreHeadings

\chapter[Extrema y media razón holográfica]{Extrema y media razón
holográfica: completación nonádica y conservación evolutiva de la unidad}
\label{chap:completacion-unidad-rev7}
\HMTClaim{HMT-R-EMRH-COMPLETION}
\HMTClaim{HMT-R-UNIT-PROJECTIVE}
La extensión cúbica no se introduce como una ilustración auxiliar. Construye
el paso causal desde el interior triádico hasta la completación base mil:
\[
999=729+243+27=729+270,
\qquad
1000=999+1=729+271.
\]
El primer miembro es la completación perforada; el segundo añade el ápice
único. Esta lectura de la extrema y media razón holográfica expresa cómo
crecen rango, frontera e información sin consumir la unidad normalizada y
prepara los cilindros que harán posible la prolongación infinita.
\HMTDemoteHeadings
\input{sections/hmt/04_cubo_emrh}
\input{sections/hmt/04b_extension_nonadica_completacion_binomial}
\HMTRestoreHeadings

\chapter{Número HMT, cilindros y evaluación arquimediana}
\label{chap:numero-cilindros-rev7}
\HMTClaim{HMT-R-NUMBER-INVERSE-SYSTEM}
\HMTClaim{HMT-R-NUMBER-ARCHIMEDEAN-SHADOW}
\HMTDemoteHeadings
\input{sections/hmt/05_tpk_numero_medicion}
\input{sections/hmt/05b_extensiones_dinamica_cociente}
\input{sections/hmt/delta_297/14c_taxonomia_monodromia_nonadica}
\HMTRestoreHeadings

\chapter[Ley de las Nueve Puertas y generación infinita]{Ley de las Nueve
Puertas: filtración coinductiva de supervivencia con monodromía nonádica
\texorpdfstring{\(\mathcal G_9\)}{G9} y generación infinita de
\texorpdfstring{\(\pi,e,\varphi\)}{pi, e y phi}}
\label{chap:generacion-constantes-rev7}
\HMTClaim{HMT-R-GEN-PI-E-PHI}
\HMTClaim{HMT-R-MONODROMY-MEMORY}
\HMTClaim{HMT-R-PI-FIVE-REGIONS}
La dinámica finita APP--TRIT--TPK no tiene por resultado una lista finita de
cifras. Construye coinductivamente una sección compatible
\[
x_\chi=(x_m)_{m\geq0}\in
\varprojlim_m\operatorname{Surv}^{\chi}_m,
\qquad \chi\in\{\pi,e,\varphi\},
\]
cuyos cilindros anidados no vacíos tienen diámetro tendente a cero. Las mil
cifras certificadas son un testigo subordinado; cualquier observación finita
posterior es un truncamiento de esta misma rama infinita.
\HMTDemoteHeadings
\chapter{Prueba finita de las regiones de
\texorpdfstring{\(\pi\), \(e\) y \(\varphi\)}{pi, e y phi}}
\label{chap:regiones}
\index{regiones de pi@regiones de \(\pi\)}

La generación precedente ha fijado qué hacen clausura, propagación
y autoescala. Este bloque contiene la prueba finita que sostiene esa
narración: catálogo de regiones, multiplicidad, orientación, transporte y
prolongación. La evaluación arquimediana se aplica al final sobre el
catálogo ya construido y no reemplaza su genealogía.

\section{Proyección ternaria del catálogo}

Sea
\[
 T_{\rm cat}:\Omega_{\rm TPK}\longrightarrow\{0,\ldots,999\}^6
\]
la aplicación de catálogo construida en el \cref{chap:tpk}, y sea
\[
 \mathcal U_6:=\operatorname{im}(T_{\rm cat}).
\]
El censo exhaustivo da \(|\mathcal U_6|=468\). Para
\(U_6=(u_1,\ldots,u_6)\in\mathcal U_6\), definimos
\[
\Psi:\mathcal U_6\longrightarrow\mathbb F_3^6,
\qquad
\Psi(U_6)=((-u_1)\bmod3,\ldots,(-u_6)\bmod3).
\]
La selección interna del estado enriquecido de TPK, demostrada en el capítulo
precedente, fija las tres etiquetas
\[
w_\pi=010211,
\qquad
w_e=201101,
\qquad
w_\varphi=121200.
\]
Esta palabra ternaria es el valor de una proyección definida sobre el
catálogo, no sobre todo el cubo entero \(\{0,\ldots,999\}^6\). Esta
precisión de dominio es necesaria: la preimagen de una palabra en el cubo
entero ambiente sería enorme,
mientras que las fibras empleadas a continuación son fibras relativas a
\(\mathcal U_6\). A cada elemento de la fibra el catálogo adjunta
multiplicidad, firma y estado de retorno.

\section{La fibra de \texorpdfstring{\(\pi\)}{pi}: cinco regiones distintas}

La microfibra relativa al catálogo
\[
\mathscr R_\pi^{\rm micro}
=\{U\in\mathcal U_6:\Psi(U)=010211\}
\]
contiene cinco regiones \(U_6\) distintas:
\[
\begin{aligned}
U_\pi^\perp&=(501,614,498,169,272,272),\\
U_{\pi,1}^{++}&=(810,923,870,169,272,272),\\
U_{\pi,2}^{++}&=(810,983,810,169,272,272),\\
U_{\pi,3}^{++}&=(870,923,810,169,272,272),\\
U_\pi^{--}&=(870,923,810,418,674,521).
\end{aligned}
\]
Sus firmas multiplicativas son, respectivamente,
\[
555555,\quad339555,\quad393555,\quad933555,\quad933717,
\]
y sus multiplicidades
\[
432,\quad144,\quad144,\quad144,\quad144.
\]

\begin{figure}[p]
\centering
\resizebox{.96\textwidth}{!}{\input{figures/fig_cinco_regiones_pi}}
\caption{Microfibra de \(\pi\). La proyección a \(w_\pi\) conserva la
palabra y elimina la diferencia regional.}
\label{fig:cinco-pi}
\end{figure}

Bajo el alineamiento regional declarado con el sistema de incidencia, las
etiquetas de rol se distribuyen como
\[
5=1_\perp+3_{++}+1_{--};
\]
por multiplicidad,
\[
1008=432+4\cdot144.
\]
Las cuatro primeras regiones conservan la cola \(169|272|272\); la última es
la rama de retorno mutado. Los identificadores de catálogo no son canónicos,
porque cambian al reordenar las filas. Los vectores \(U_6\) sí lo son dentro
del catálogo fijado.

\begin{remark}[Fibra regional y cociente ternario]
El cociente
\(\mathscr R_\pi^{w_6}=\{010211\}\) tiene un elemento porque elimina los datos
regionales. La microfibra tiene cinco porque los conserva. Escribir «la región
de \(\pi\)» sólo es admisible si se especifica cuál de los dos tipos se está
usando.
\end{remark}

\section{Las fibras simples de \texorpdfstring{\(e\)}{e} y
\texorpdfstring{\(\varphi\)}{phi}}

La misma proyección que produce la microfibra quíntuple de \(\pi\)
determina dos fibras simples:
\[
 \Psi^{-1}(201101)=\{U_e\},
 \qquad
 \Psi^{-1}(121200)=\{U_\varphi\},
\]
donde
\[
 U_e=(850,963,890,149,252,212),
 \qquad
 U_\varphi=(830,943,830,109,252,252).
\]
Ambas regiones tienen multiplicidad \(144\), clase de orientación aditiva
\(\mathrm{ES}\), soporte multiplicativo de cardinal seis y tipo residual
\(A\). Sus firmas multiplicativas
son, respectivamente,
\[
 e_\times(U_e)=771339,
 \qquad
 e_\times(U_\varphi)=555933.
\]
La reducción directa de sus coordenadas es
\[
 U_e\bmod3=102202,
 \qquad
 U_\varphi\bmod3=212100;
\]
por la convención \(\Psi(U)=(-U)\bmod3\), estas dos palabras se convierten
en \(201101\) y \(121200\). Así, la diferencia entre la fibra de \(\pi\) y
las de \(e,\varphi\) ya está presente antes de toda evaluación decimal:
\(\pi\) posee cinco realizaciones regionales con multiplicidad total
\(1008\), mientras que \(e\) y \(\varphi\) poseen una realización cada una.

\section{Caracteres estructurales de las tres fibras}

Los datos anteriores distinguen tres comportamientos combinatorios. La fibra
de \(\pi\) registra clausura y orientación porque contiene ramas de retorno
diferentes; la de \(e\) queda caracterizada, después de aplicar la carta
arquimediana, por una complementariedad nonádica; y la de \(\varphi\) se
compara con el operador clásico de autoescala. Esta lectura se resume como
\[
\begin{array}{ccl}
\pi&:&\text{clausura, borde y orientación},\\
e&:&\text{propagación y crecimiento},\\
\varphi&:&\text{autoescala y proporción},\\
\alpha&:&\text{acoplamiento de las tres evaluaciones mediante el vector }K.
\end{array}
\]
La primera asignación es un invariante del catálogo regional. Las otras dos
utilizan los operadores arquimedianos que se introducen en las secciones
siguientes: la complementariedad de \(e\) y el radio espectral de
\(\varphi\). El cierre dodecafásico construye después el vector \(K\), la
frontera y la identidad conjunta que genera \(\alpha\); esta prueba finita no
introduce una calibración paralela de esos objetos.

Las doce tríadas arquimedianas que alimentan el cierre dodecafásico son
\[
\begin{aligned}
P={}&(141,592,653,589,793,238,462,643,383,279,502,884),\\
E={}&(718,281,828,459,045,235,360,287,471,352,662,497),\\
\Phi={}&(618,033,988,749,894,848,204,586,834,365,638,117).
\end{aligned}
\]
Estos tres vectores no son salidas de \(\Psi\): pertenecen al segundo nivel
de la construcción. El control finito que se especifica más abajo reproduce
sus cuatro primeras tríadas. Las doce tríadas completas son la salida de la
generación monodrómica del capítulo precedente y no se introducen como datos
de este control de cuatro pasos. En particular, \(P\) es la coordenada
decimal común de las cinco regiones de \(\pi\), pero no las identifica en el
espacio completo de estados.

\section{La región de \texorpdfstring{\(e\)}{e} y la complementariedad \texorpdfstring{\(1000-1\)}{1000-1}}

La primera ventana arquimediana asignada a \(w_e\) es
\[
718|281|828|459.
\]
Entre las \(243\) palabras de
\(\mathcal W_6:=\Psi(\mathcal U_6)\), y después de aplicar el control finito
de la carta arquimediana definido en la sección siguiente, \(w_e\) es la única
cuya primera pareja de tríadas suma \(999\). Además,
\[
718=729-11,
\qquad
281=271+10,
\qquad
718+281=999=1000-1,
\]
\[
828=3\cdot271+15,
\qquad
459=\sum R^\times.
\]
La caracterización finita es exacta una vez fijados el catálogo, los
operadores del TPK y la condición de complementariedad. Al aplicar el
predicado no se introduce otro valor decimal; esta unicidad es un control
interno de la generación ya construida y no una segunda derivación de \(e\).

El objeto examinado en esta sección es la órbita finita de cuatro
transiciones. La extensión a profundidad arbitraria y su compatibilidad con
la supervivencia no se infieren de esta ventana: pertenecen a la relación
EF--\(G_9\) y al teorema monodrómico del capítulo precedente.

\section{\texorpdfstring{\(\varphi\)}{phi}, Fibonacci y nivel cinco}

La proporción áurea es el radio espectral de
\[
N_\tau=\begin{pmatrix}0&1\\1&1\end{pmatrix},
\]
y el cociente continuo asociado satisface el punto fijo
\[
\varphi=1+\frac1\varphi.
\]
Estas son identidades clásicas. Una extensión posible considera el nivel
cinco de la correspondencia pentada--hexada, las clases
de Rogers--Ramanujan y las reglas de Fibonacci. Para convertir la convergencia
\(R(q)\to\varphi^{-1}\) en información HMT no basta evaluar cerca de
\(q=1\), porque el límite es universal. El observable debe ser diferencial o
espectral después de sustraer la asintótica dominante. La construcción de tal
observable y de su entrelazador con los productos modulares constituye una
hipótesis adicional a los resultados de esta sección.

\begin{remark}[Correspondencia pentada--hexada]
Una ruta concreta consiste en seleccionar mediante HMT un elemento de orden
cinco en la acción Witt, descomponer sus modos en las clases
\(\{\pm1\}\) y \(\{\pm2\}\pmod5\), y demostrar que las funciones
generatrices resultantes son los dos productos de Rogers--Ramanujan. Si el
operador de transferencia es \(N_\tau\), \(\varphi\) aparecerá como radio
espectral y la relación con anyones de Fibonacci será una representación, no
una analogía verbal.
\end{remark}

\section{Verificación finita de la contracción arquimediana sobre cilindros regionales}

El orden de tipos es
\[
\begin{aligned}
\text{admisibilidad}
&\to\text{semillas}
\to\text{palabra corta}
\to\text{calendario de continuación}\\
&\to\text{cilindro}
\to\text{supervivencia}
\to\text{evaluación real}.
\end{aligned}
\]
La cadena decimal aparece al final. Para hacer explícita la evaluación
finita, fijamos sobre \(\mathbb F_3^6\), con vectores fila y acción por la
derecha, las matrices
\[
L_0=
\begin{pmatrix}
2&2&2&1&2&1\\
2&1&2&2&1&1\\
1&0&0&1&0&2\\
0&0&1&1&1&0\\
2&2&2&0&2&1\\
2&1&2&1&0&0
\end{pmatrix},
\qquad
L_1=
\begin{pmatrix}
2&1&2&0&2&2\\
1&2&1&1&2&0\\
1&2&1&1&1&2\\
0&1&0&0&2&1\\
2&0&2&1&0&1\\
0&2&2&2&1&1
\end{pmatrix}.
\]
Estas matrices son operadores del estado enriquecido de TPK. No son parámetros que
se ajusten durante la evaluación: fijados el calendario y el estado inicial,
la lectura es determinista. Además, las seis transiciones publicadas de rango
completo caracterizan cada operador de manera única,
\[
L=X^{-1}Y\qquad\text{en }\mathbb F_3.
\]
Esta caracterización dual pertenece al estado enriquecido de TPK. Obtener de nuevo todos
sus campos desde la proyección reducida APP mínima sería un teorema de
realización distinto y no convierte \(L_0,L_1\) en datos introducidos a mano.

Para \(w=(d_0,\ldots,d_5)\in\mathbb F_3^6\) y un calendario
\(c=(c_1,c_2,c_3,c_4)\in\{0,1\}^4\), sean
\[
 u_0=w,
 \qquad
 u_j=u_{j-1}L_{c_j}\quad(1\leq j\leq4),
 \qquad
 W_c(w)=u_0\Vert u_1\Vert u_2\Vert u_3\Vert u_4.
\]
Si \(W_c(w)=(a_0,\ldots,a_{29})\), definimos mediante aritmética entera
\[
 N_c(w)=\sum_{k=0}^{29}a_k3^{29-k},
 \qquad
 \mathcal A_c(w)_j=
 \left\lfloor\frac{1000^jN_c(w)}{3^{30}}\right\rfloor\bmod1000,
 \quad1\leq j\leq4.
\]
La aplicación
\(\mathcal A_{0011}:\mathbb F_3^6\to\{0,\ldots,999\}^4\) es, por tanto,
la ventana finita de la carta arquimediana. Su dominio y codominio difieren de
los de \(\Psi\), y su definición abarca cuatro iteraciones. La prolongación a
profundidad arbitraria no se postula a partir de esta ventana: la proporciona
la conexión EF--\(G_9\) ya construida.

El censo exacto de las \(729\) palabras verifica que \(\mathcal A_{0011}\)
es inyectiva y que
\[
\begin{array}{c|c}
w&\mathcal A_{0011}(w)\\ \hline
010211&(141,592,653,589)\\
201101&(718,281,828,459)\\
121200&(618,033,988,749)
\end{array}
\]
son reconocimientos únicos. Entre los dieciséis calendarios binarios,
\(0011\) es el único que conserva simultáneamente estos tres
reconocimientos. Las \(144\) perturbaciones obtenidas al sustituir una entrada
de \(L_0\) o \(L_1\) por uno de los otros dos residuos destruyen el triple
reconocimiento. Los órdenes son
\[
\operatorname{ord}(L_0)=80,
\quad
\operatorname{ord}(L_1)=26,
\quad
\operatorname{ord}(L_0L_1)=242.
\]

La misma comprobación algebraica fija el estatuto lógico de este resultado.
Para \(L_0\), los seis vectores de entrada que proceden de las dos
primeras transiciones de \(w_\pi,w_e,w_\varphi\) tienen rango seis sobre
\(\mathbb F_3\); para \(L_1\) ocurre lo mismo con los seis vectores de las dos
transiciones siguientes. Por ello, las transiciones determinan cada matriz de
manera única. La unicidad, los órdenes y la rigidez frente a perturbaciones
son propiedades exactas del operador fijado. El cálculo directo distingue
esta caracterización dentro del estado completo de una eventual
reconstrucción adicional desde APP mínima.

La proyección observable de APP y el espacio completo de estados se distinguen
mediante la proyección de estado definida en el \cref{chap:tpk}; la carta
arquimediana es una aplicación posterior sobre el cociente ternario.

\begin{remark}[Catálogo y evaluación]
El resultado finito probado en esta sección tiene la forma
\[
\Omega_{\rm TPK}\xrightarrow{T_{\rm cat}}\mathcal U_6
\xrightarrow{\Psi}\mathcal W_6
\xrightarrow{\mathcal A_{0011}}\{0,\ldots,999\}^4.
\]
Las dos primeras flechas constituyen el catálogo; la tercera es el control
finito de la carta. Conservar esta separación impide atribuir al censo regional una
afirmación que sólo se ha verificado después de fijar \(L_0,L_1\).
\end{remark}

\input{sections/hmt/06b_cinco_ciclos_pi}
\input{sections/hmt/06c_canales_enteros_constantes}
\input{sections/hmt/06d_biografias_numeros_rev6}
\input{sections/hmt/06e_certificado_generacion_monodromica_rev7}
\input{sections/hmt/14_numero_heisenberg_y_d108}
\HMTRestoreHeadings

\chapter{Cierre dodecafásico de
\texorpdfstring{\(\alpha_{\mathrm{HMT}}\)}{alpha HMT}}
\label{chap:alpha-rev7}
\HMTClaim{HMT-R-ALPHA-DODECAPHASE}
\HMTClaim{HMT-R-DODEC-EXTRACTOR}
El cierre de \(\alpha_{\rm HMT}\) recibe dos lecturas coordinadas del mismo
estado terminal enriquecido. No se infiere la coordenada firmada a partir de
los tres bloques visibles: ambos son proyecciones distintas del estado TPK.
\HMTDemoteHeadings
\input{sections/hmt/08_dodecafase_alpha}
\input{sections/hmt/08a_alpha_dos_vias_t1_rev8}
\HMTRestoreHeadings

\chapter{Conúcleo determinantal y orden decimal de la acción}
\label{chap:escala-accion-rev7}
\HMTClaim{HMT-R-ACTION-COKERNEL}
\HMTClaim{HMT-R-ACTION-MINUS34}
La forma normal de Smith del lector local fija el índice dieciséis y, con él,
\(\Sigma_H=\varphi\alpha_{\rm HMT}^{16}\) en la década \(-34\) antes de
elegir la carta \(\mathrm{J\,s}\). La escala de acción no se importa desde
Planck: la metrología sólo reconoce después la coordenada ya construida.
\HMTDemoteHeadings
\input{sections/hmt/delta_297/16g_escala_accion_determinantal}
\HMTRestoreHeadings

\chapter{Bisagra dodecafásica--Witt}
\label{chap:entrelazador-dodecafase-witt-rev8}
La rama excepcional recibe aquí su entrelazador propio
\(P_G\leftrightarrow P_3\leftrightarrow Q_3\). Conserva rango,
orientación y medida, pero no genera la salida arquimediana ni sustituye al
extractor dodecafásico.
\HMTDemoteHeadings
\input{sections/hmt/delta_297/16c_entrelazador_dodecafase_witt}
\HMTRestoreHeadings

\chapter[Rama arquimediana post-\texorpdfstring{\(\alpha\)}{alpha}]{Rama
arquimediana post-\texorpdfstring{\(\alpha\)}{alpha}: transportador positivo
y coordenadas angulares internas}
\label{chap:entrelazador-angular-rev8}
\HMTClaim{HMT-R-CBAND}
Fijados el estado enriquecido común y el cierre dodecafásico, esta rama
desarrolla la salida arquimediana: acción positiva, ángulo, contraángulo y
canales internos. Esta residencia evita que Barbero, el residuo angular o la
ley de masas parezcan generar retroactivamente las coordenadas \(A,C^*\).
\HMTDemoteHeadings
\input{sections/hmt/16c_transportador_positivo_accion}
\input{sections/hmt/delta_297/16d_medida_accion_positiva}
\input{sections/hmt/delta_297/16e_operador_angular_y_dictamen}
\input{sections/hmt/delta_297/16f_realizacion_analitica_contraangulo}
\input{sections/hmt/delta_297/16h_invariante_focal_relativo_rev8}
\HMTRestoreHeadings

\chapter{Raíz interna de \texorpdfstring{\(-1\)}{-1} y extensión
exponencial de Euler}
\label{chap:raiz-euler-rev7}
\HMTClaim{HMT-R-SQRT-MINUS-ONE}
\HMTClaim{HMT-R-EULER-EXTENSION}
\HMTDemoteHeadings
\chapter{Construcción de Paley, estructura compleja finita y código ternario
extendido de Golay}
\label{chap:paley-codigo-witt}
\index{Paley!matriz de conferencia}
\index{Golay!código ternario extendido}
\index{Witt!sistema de Steiner}

La prolongación excepcional de la frontera nonádica requiere una construcción
algebraica anterior a cualquier uso de hexadas, octadas o grupos de Mathieu.
Esta sección fija esa construcción en coordenadas. El orden lógico es
\[
 (\mathbb Z/9\mathbb Z)^\times
 \longrightarrow \mathbb P^1(\mathbb F_5)
 \longrightarrow C_P
 \longrightarrow A_W
 \longrightarrow C_W
 \longrightarrow W_{12}.
\]
Las flechas no se deducen de coincidencias numéricas: las dos primeras fijan
un sistema de coordenadas, las dos siguientes son operaciones matriciales y la
última toma soportes proyectivos de peso seis.

\section{Coordenadas proyectivas y matriz de conferencia}

El conjunto de unidades módulo nueve tiene seis elementos. Elegimos la
bijección de coordenadas
\[
\begin{array}{c|cccccc}
u&1&2&4&8&7&5\\ \hline
\theta(u)&\infty&0&1&2&3&4
\end{array}
\]
en la que la primera fila sigue las potencias del generador \(2\) módulo
nueve. Esta elección fija un origen y una orientación, e identifica el
conjunto subyacente con la recta proyectiva
\[
 \mathbb P^1(\mathbb F_5)=\{\infty,0,1,2,3,4\}.
\]
La aplicación \(\theta\) es un calibre de coordenadas, no un homomorfismo ni
una consecuencia canónica de la estructura multiplicativa de las unidades.
El carácter cuadrático de \(\mathbb F_5\) determina, en el orden anterior, la
matriz simétrica de conferencia
\begin{equation}
 C_P=
 \begin{pmatrix}
 0& 1& 1& 1& 1& 1\\
 1& 0& 1&-1&-1& 1\\
 1& 1& 0& 1&-1&-1\\
 1&-1& 1& 0& 1&-1\\
 1&-1&-1& 1& 0& 1\\
 1& 1&-1&-1& 1& 0
 \end{pmatrix}.
 \label{eq:paley-conference-explicita}
\end{equation}
El producto fila por fila da
\begin{equation}
 C_PC_P^{\mathsf T}=5I_6.
 \label{eq:paley-conference-identidad}
\end{equation}
La diagonal nula y la ortogonalidad de filas expresan, respectivamente, la
ausencia de autoincidencia y el equilibrio del carácter cuadrático. Un cambio
simultáneo de orden o de signos produce una matriz monomialmente equivalente;
la matriz \eqref{eq:paley-conference-explicita} fija el calibre empleado en
todo el manuscrito.

\section[Reducción ternaria y estructura compleja interna]{Reducción
ternaria y endomorfismo cuadrático
\texorpdfstring{\(A_W^2=-I_6\)}{AW2=-I6}}

Reducimos las entradas de \(C_P\) módulo tres, identificando
\(-1\equiv2\pmod3\). Se obtiene
\begin{equation}
 A_W=
 \begin{pmatrix}
 0&1&1&1&1&1\\
 1&0&1&2&2&1\\
 1&1&0&1&2&2\\
 1&2&1&0&1&2\\
 1&2&2&1&0&1\\
 1&1&2&2&1&0
 \end{pmatrix}
 \in M_6(\mathbb F_3).
 \label{eq:AW-explicita}
\end{equation}
Como \(A_W=A_W^{\mathsf T}\), la reducción de
\eqref{eq:paley-conference-identidad} implica
\begin{equation}
 A_W^2=A_WA_W^{\mathsf T}=2I_6=-I_6.
 \label{eq:AW-raiz-menos-uno}
\end{equation}
Por tanto \(A_W\) realiza una raíz cuadrada de \(-I\) dentro de un espacio
vectorial definido sobre \(\mathbb F_3\). No se identifica el cuerpo ternario
con \(\mathbb C\); se afirma la identidad matricial exacta
\eqref{eq:AW-raiz-menos-uno}, que proporciona una estructura compleja finita
en el sentido algebraico de un endomorfismo \(J\) con \(J^2=-I\).

\begin{proposition}[Descomposición en subespacios propios conjugados]
Sea \(V=\mathbb F_3^6\) y sea \(J=A_W\). La extensión escalar
\(V\otimes_{\mathbb F_3}\mathbb F_9\) se descompone en los dos espacios
propios de \(J\) correspondientes a las raíces \(\iota\) y \(-\iota\) de
\(X^2+1\) en \(\mathbb F_9\).
\end{proposition}

\begin{proof}
El polinomio \(X^2+1\) no tiene raíces en \(\mathbb F_3\), pues los
cuadrados son cero y uno. En \(\mathbb F_9\) se factoriza como
\((X-\iota)(X+\iota)\), con raíces distintas. La identidad \(J^2=-I\)
muestra que el polinomio mínimo de \(J\) divide a \(X^2+1\); como \(J\) no
es escalar, coincide con él. El endomorfismo es, por tanto, diagonalizable
después de extender escala a \(\mathbb F_9\), y la suma de ambos espacios
propios es directa y exhaustiva.
\end{proof}

\section{Código gráfico autodual}

Definimos el morfismo lineal
\[
 c:\mathbb F_3^6\longrightarrow\mathbb F_3^{12},
 \qquad c(q)=(q,qA_W),
\]
y su imagen
\[
 C_W=\operatorname{im}c.
\]
La matriz generadora es \(G=(I_6\mid A_W)\). Por
\eqref{eq:AW-raiz-menos-uno},
\[
 GG^{\mathsf T}=I_6+A_WA_W^{\mathsf T}=I_6+2I_6=0
 \quad\text{en }\mathbb F_3.
\]
Las seis filas de \(G\) son independientes, de modo que \(C_W\) tiene
dimensión seis y es igual a su dual ortogonal. La enumeración de los
\(3^6=729\) vectores \(q\) proporciona el enumerador de pesos
\begin{equation}
 W_{C_W}(y)=1+264y^6+440y^9+24y^{12}.
 \label{eq:enumerador-CW}
\end{equation}
En particular, la distancia mínima es seis y
\[
 C_W=[12,6,6]_3,
\]
el código ternario extendido de Golay en las coordenadas fijadas por
\eqref{eq:AW-explicita}.

La enumeración exhaustiva de los \(729\) mensajes, formando cada palabra con
la matriz impresa y contando sus pesos, produce
\eqref{eq:enumerador-CW}.

\begin{proposition}[Palabra marcada y controles uniformes]
\label{prop:palabra-marcada-witt}
En la presentación orientada determinada por el corte \(6+6\), la matriz
\(A_W\) de \eqref{eq:AW-explicita} y la orientación positiva, sea
\[
q_\ast=(1,1,1,1,1,1)\in\mathbb F_3^6.
\]
Entonces
\[
\boxed{
c_\ast=c(q_\ast)
=(1,1,1,1,1,1,2,1,1,1,1,1)\in C_W.}
\]
En cambio,
\[
(1,\ldots,1)\notin C_W,
\qquad
(2,\ldots,2)\notin C_W.
\]
La primera palabra es canónica relativamente a la presentación orientada
declarada; no es una palabra absoluta independiente del calibre.
\end{proposition}

\begin{proof}
La suma de las filas de la matriz impresa da
\[
q_\ast A_W=(2,1,1,1,1,1),
\]
que prueba la fórmula de \(c_\ast\). Para las dos palabras uniformes, la
primera mitad fija respectivamente \(q=q_\ast\) y \(q=-q_\ast\), mientras
las segundas mitades requeridas serían \(q_\ast\) y \(-q_\ast\); las
segundas mitades efectivas son
\((2,1,1,1,1,1)\) y \((1,2,2,2,2,2)\). Por tanto ninguna palabra uniforme
pertenece al código. La enumeración exhaustiva de las \(729\) palabras
confirma que no aparece otra excepción.
\end{proof}

\[
\mathsf R=
\begin{pmatrix}
0&1\\
1&0
\end{pmatrix},
\qquad
\mathsf R C\mathsf R=C^{-1}.
\]
La palabra marcada especializa, después de la construcción abstracta del
\cref{thm:funtor-ciclotomico-fase-longitud}, la familia de trivializaciones
mediante
\[
c_b:=c_{b+1},
\qquad c\in C_W,\quad b\in\{0,\ldots,11\},
\]
donde la expresión de la derecha usa la numeración impresa
\(1,\ldots,12\). Esta convención se mantiene en todo lo que sigue y tipa, en
particular, \((c_\ast)_b\) y el posterior \(P_b(c)\). Entonces
\[
P_b^{(\ast)}=
\begin{cases}
\mathsf R,&(c_\ast)_b=1,\\
I,&(c_\ast)_b=2.
\end{cases}
\]
Esta regla depende del corte, de la matriz y de la orientación declarados;
no interviene en la definición previa del funtor fase--longitud.

La ausencia de las palabras uniformes es un control de esta orientación:
impide usarlas en la prueba concreta del vecino, pero no altera por sí sola
el pegado Niemeier, que se conserva por la acción trivial de cada Coxeter
sobre \(A_2^\ast/A_2\).

\section{Soportes proyectivos y sistema de Witt}

Cada una de las 264 palabras de peso seis tiene un soporte de seis
coordenadas. Las palabras opuestas \(u\) y \(-u\) poseen el mismo soporte y
no hay otra identificación, por lo que se obtienen
\[
 \frac{264}{2}=132
\]
hexadas distintas. Denotemos por \(\mathcal H_W\) la familia de esos
soportes.

\begin{theorem}[Sistema de Witt ternario]
El sistema de incidencia
\[
 \bigl(\{1,\ldots,12\},\mathcal H_W\bigr)
\]
es un sistema de Steiner \(S(5,6,12)\). En consecuencia, es una realización
del diseño de Witt \(W_{12}\).
\end{theorem}

\begin{proof}
El recorrido de los \(\binom{12}{5}=792\) subconjuntos de cinco puntos cuenta
las hexadas que contienen cada uno. Todos los conteos valen
uno. Ésta es exactamente la propiedad definitoria de un
\(S(5,6,12)\). El número de bloques también resulta de la identidad de
incidencias
\[
 |\mathcal H_W|\binom65=\binom{12}{5},
\]
que da \(|\mathcal H_W|=132\).
\end{proof}

Esta sección fija por tanto, antes de utilizarla, la cadena completa
\[
 C_P\longmapsto A_W\longmapsto C_W\longmapsto W_{12}.
\]
Las secciones posteriores añaden el marco dodecafásico, las involuciones de
estrella, la elevación relativa a una dodecada binaria y la construcción
reticular. Ninguno de esos pasos modifica el calibre matricial ya establecido.

\input{sections/hmt/09b_algebras_cuadraticas}
\HMTRestoreHeadings

\chapter{Incidencia ternaria, las veintisiete rectas y
\texorpdfstring{\(W(E_6)\)}{W(E6)}}
\label{chap:incidencia-e6-rev7}
\HMTClaim{HMT-R-LINES27-E6}
\HMTClaim{HMT-R-MINUS-ONE-TWELFTH}
\HMTDemoteHeadings
\input{sections/hmt/07_nucleo_hensel_witt}
\input{sections/hmt/15b_holonomia_schlafli_e6}
\HMTRestoreHeadings

\chapter{Dominio enriquecido común y doble evaluación global}
\label{chap:doble-evaluacion-rev7}
\HMTClaim{HMT-R-DOUBLE-PROJECTION-COMMON-DOMAIN}
Un mismo estado enriquecido admite dos evaluaciones HMT, no tres lectores
paralelos: la rama arquimediana contrae cilindros y la rama excepcional
conserva incidencia, gradación, orientación y acarreo. La prolongación de
Mecánica Dimensional comenzará después, conservando la genealogía de ese
mismo antecedente.
\HMTDemoteHeadings
\input{sections/hmt/09g_lector_excepcional_global_rev6}
\input{sections/hmt/09_doble_proyeccion}
\input{sections/hmt/delta_297/09e_globalizacion_doble_evaluacion}
\input{sections/hmt/delta_297/09c_correspondencia_armonica_finita}
\input{sections/hmt/delta_297/09f_obstruccion_lector_excepcional_minimo}
\HMTRestoreHeadings

\chapter{Corredor excepcional ramificado}
\label{chap:corredor-excepcional-rev7}
\HMTClaim{HMT-R-EXCEPTIONAL-RAMIFIED}
\HMTClaim{HMT-R-MONSTER-MOONSHINE}
La diferencia integrada \(459-405=54\) reaparece como coeficiente de la traza
graduada, pero no se convierte por semejanza en un objeto del Monstruo. El
corredor preserva sus mapas: desde \(C_W\) se separan la rama
\(W_{12}\to W_{24}\) y el pegado reticular
\(N(A_2^{12})\to\Lambda_{24}\); sólo después aparecen
\(V_\Lambda\), \(V^\natural\), el Monstruo y Moonshine.
\HMTDemoteHeadings
\input{sections/hmt/03f_acoplamiento_app_fock}
\input{sections/hmt/15_bandera_y_rama_excepcional}
\input{sections/hmt/15c_cierre_micro_macro_nivel_tres}
\HMTRestoreHeadings

\chapter{Reducciones de rango y dualidad residuo--cociente}
\label{chap:pantallas-dualidad-rev7}
\HMTClaim{HMT-R-DIMENSIONAL-SCREENS}
\HMTClaim{HMT-R-T-DUALITY}
\HMTClaim{HMT-R-M-THEORY-MATH}
\HMTDemoteHeadings
\input{sections/hmt/16b_pantallas_dualidad_teoria_m}
\input{sections/hmt/16d_pantalla_cubica_firma_31}
\HMTRestoreHeadings

\chapter{Aritmética nativa, vacancias y desarrollos espectrales}
\label{chap:aritmetica-espectral-rev7}
\HMTClaim{HMT-R-NATIVE-PRODUCT}
\HMTClaim{HMT-R-VACANCIES}
\HMTClaim{HMT-R-PELL}
\HMTClaim{HMT-R-PRIME-CLOCKS}
\HMTClaim{HMT-R-FEIGENBAUM-FINITE}
\HMTDemoteHeadings
\input{sections/hmt/delta_297/06c_espectro_lector_243}
\input{sections/hmt/delta_297/14b_cobertura_y_taxonomia_numeros}
\input{sections/hmt/10_producto_nativo}
\input{sections/hmt/18_funcionales_espectrales_triadicos}
\input{sections/hmt/19_rotaciones_vacancias_criticidad}
\input{sections/hmt/19a_constantes_torre_triadica_rev6}
\HMTRestoreHeadings

\chapter{Continuo, memoria e incidencia}
\label{chap:sintesis-hmt-rev7}
\HMTClaim{HMT-R-HMT-SYNTHESIS}
\HMTDemoteHeadings
\input{sections/hmt/17_sintesis_rev6}
\HMTRestoreHeadings

\part{Mecánica Dimensional}

\chapter{Principio de ambivalencia, persistencia de extensiones y definición
de partícula}
\label{chap:particula-rev7}
\HMTClaim{MD-R-AMBIVALENCE}
\HMTClaim{MD-R-PARTICLE-CORE}
\HMTClaim{MD-R-VIRTUAL-PARTICLE}
\HMTClaim{MD-R-ROUTE-LEDGER-SIGNATURE}
Mecánica Dimensional se abre con el estado completo, no con una magnitud ya
proyectada. El principio de ambivalencia distingue las lecturas areal y
volumétrica mientras la memoria conserva su origen común. Sobre esa base una
partícula se define como sección persistente de una extensión, y su ruta se
lee antes de asignarle masa, carga o nombre físico.
\HMTDemoteHeadings
\input{sections/md/01_medida_ambivalencia}
\def\HMTIncludeParticleCore{1}
\input{sections/md/04_particula_estadistica}
\let\HMTIncludeParticleCore\undefined
\chapter{Banda de partícula, anti-sección y virtualidad}
\label{chap:banda-particula-virtual}
\index{partícula!banda de ruta}
\index{antipartícula!anti-sección orientada}
\index{partícula virtual}

La persistencia de una ruta produce una partícula antes de que se evalúe su
masa. Esta prioridad permite distinguir tres objetos:

\begin{enumerate}
\item la \emph{sección orientada}, que conserva la dirección de la ruta;
\item la \emph{banda}, que reúne una sección y su conjugada;
\item la \emph{persistencia}, que decide si la firma alcanza o no una rama
global.
\end{enumerate}

\begin{definition}[Palabras dodecafásicas y conjugación material]
\label{def:palabras-cm}
Sea
\[
 \mathscr T_{12}^{\rm word}:=\{-1,0,1\}^{12}.
\]
Para \(\mathbf t=(t_1,\ldots,t_{12})\) escribimos
\[
 \operatorname{rev}(\mathbf t)
 =(t_{12},\ldots,t_1),
 \qquad
 C_M(\mathbf t):=-\operatorname{rev}(\mathbf t).
\]
La inversión \(C_M\) cambia simultáneamente el sentido de lectura y la
orientación trítica.
\end{definition}

\begin{theorem}[Involución material y sector fijo]
\label{thm:involucion-cm-censo}
La aplicación \(C_M\) es una involución de
\(\mathscr T_{12}^{\rm word}\) y
\[
 \left|\operatorname{Fix}(C_M)\right|=3^6=729.
\]
\end{theorem}

\begin{proof}
La reversión y el cambio de signo conmutan y ambos tienen cuadrado igual a
la identidad; por tanto \(C_M^2=I\). Una palabra es fija si y sólo si
\[
 t_i=-t_{13-i},\qquad 1\leq i\leq6.
\]
Las seis primeras coordenadas son libres y las seis últimas quedan
determinadas por ellas. Cada coordenada libre admite tres valores, de modo
que el sector fijo tiene cardinal \(3^6\).
\end{proof}

Esta conjugación finita no se obtiene cambiando únicamente el signo visible:
invierte también la orientación del registro.

\section{La banda de ruta}

Sea \(\mathscr T_{12}^{\rm word}\) el espacio de palabras de la
\cref{def:palabras-cm}. Añadimos una orientación de sección
\(\chi\in\{-1,+1\}\) y definimos
\[
(\mathbf t,\chi)\sim(C_M\mathbf t,-\chi).
\]

\begin{definition}[Banda dodecafásica de ruta]
\label{def:banda-dodecafasica-ruta}
La banda de la sección orientada \((\mathbf t,\chi)\) es la clase
\[
[(\mathbf t,\chi)]_M
=\{(\mathbf t,\chi),(C_M\mathbf t,-\chi)\},
\]
y el espacio de bandas es
\[
\mathcal B_{12}
:=
\bigl(\mathscr T_{12}^{\rm word}\times\{-1,+1\}\bigr)/{\sim}.
\]
\end{definition}

\begin{proposition}[Libertad de la conjugación orientada]
\label{prop:accion-cm-orientada-libre}
La acción
\[
(\mathbf t,\chi)\longmapsto(C_M\mathbf t,-\chi)
\]
es libre. En consecuencia,
\[
|\mathcal B_{12}|=3^{12}.
\]
\end{proposition}

\begin{proof}
Un punto fijo exigiría \(\chi=-\chi\), imposible para
\(\chi\in\{-1,+1\}\). Todas las órbitas tienen dos elementos y el dominio
posee \(2\cdot3^{12}\).
\end{proof}

La banda no borra la orientación: la conserva como doble cubierta. La
partícula y la antipartícula son las dos secciones orientadas de la misma
banda.

\section{Masa par y carga impar}

Sea \(E_0:\mathscr T_{12}^{\rm word}\to\mathbb R\) un carácter aditivo de
ruta. Sus componentes par e impar respecto de \(C_M\) son
\[
E_+(\mathbf t)
=\frac12\bigl(E_0(\mathbf t)+E_0(C_M\mathbf t)\bigr),
\]
\[
E_-(\mathbf t)
=\frac12\bigl(E_0(\mathbf t)-E_0(C_M\mathbf t)\bigr).
\]
Entonces
\[
E_+(C_M\mathbf t)=E_+(\mathbf t),
\qquad
E_-(C_M\mathbf t)=-E_-(\mathbf t).
\]

\begin{definition}[Carácter de banda]
\label{def:caracter-banda}
Para una sección orientada definimos
\[
E_{\rm band}(\mathbf t,\chi)
=E_+(\mathbf t)+\chi E_-(\mathbf t),
\qquad
m(\mathbf t,\chi)
=m_\circ\exp\bigl(E_{\rm band}(\mathbf t,\chi)\bigr),
\]
donde \(m_\circ\) es la unidad interna de comparación de la familia.
\end{definition}

\begin{theorem}[Principio de conjugación de masa]
\label{thm:principio-conjugacion-masa-banda}
El carácter y la masa descienden al cociente de bandas:
\[
E_{\rm band}(C_M\mathbf t,-\chi)
=E_{\rm band}(\mathbf t,\chi),
\]
\[
\boxed{
m(C_M\mathbf t,-\chi)=m(\mathbf t,\chi).}
\]
Los invariantes pares pertenecen a la banda; los invariantes impares
distinguen la orientación de sus dos secciones.
\end{theorem}

\begin{proof}
Por las paridades anteriores,
\[
\begin{aligned}
E_{\rm band}(C_M\mathbf t,-\chi)
&=E_+(C_M\mathbf t)-\chi E_-(C_M\mathbf t)\\
&=E_+(\mathbf t)+\chi E_-(\mathbf t).
\end{aligned}
\]
La igualdad de masas se obtiene al aplicar la exponencial.
\end{proof}

En particular, una forma lineal con pesos simétricos es impar, mientras una
forma cuadrática compatible con la reversión es par. La carga y la
orientación cambian; la masa, construida sobre la banda, permanece.

\begin{corollary}[Electrón y positrón]
\label{cor:electron-positron-banda}
En la ecuación electrónica de banda, el par
\((n,\chi)=(-22,+1)\) y su conjugado \((+22,-1)\) producen el mismo término
\(\chi n=-22\). Por tanto, electrón y positrón poseen la misma masa y
orientaciones de carga opuestas.
\end{corollary}

\section{Cierre finito del atlas bajo anti-sección}

En las ochenta y una rutas CT--108 aparecen siete formas tríticas. El censo
exhaustivo distingue tres operaciones:
\[
\mathbf t\longmapsto-\mathbf t,\qquad
\mathbf t\longmapsto\operatorname{rev}(\mathbf t),\qquad
\mathbf t\longmapsto-\operatorname{rev}(\mathbf t).
\]

\begin{proposition}[Cierre de las formas CT--108]
\label{prop:cierre-antiseccion-ct108}
El conjunto de formas tríticas de CT--108 no es estable bajo signo puro ni
bajo reversión pura, y es estable bajo \(C_M=-\operatorname{rev}\). Sus
órbitas poseen multiplicidades
\[
\boxed{81=54+9+9+9.}
\]
La órbita de multiplicidad \(54\) es auto-conjugada; las otras tres
intercambian pares de formas.
\end{proposition}

\begin{proof}
La forma fija es
\[
\texttt{+++---+++---}
\]
y posee multiplicidad \(54\). Las tres parejas
\[
\begin{array}{c}
\texttt{+++--0+++--0}\ /\ \texttt{0++---0++---},\\
\texttt{+++-0-+++-0-}\ /\ \texttt{+0+---+0+---},\\
\texttt{+++0--+++0--}\ /\ \texttt{++0---++0---}
\end{array}
\]
poseen multiplicidad total \(9\) cada una. La aplicación
\(-\operatorname{rev}\) fija la primera y permuta los miembros de cada
pareja. El cambio de signo y la reversión por separado producen palabras
fuera de esta lista. La suma de multiplicidades es
\(54+9+9+9=81\).
\end{proof}

El resultado no identifica el sector \(54\) con otro objeto de cardinal
homónimo. Demuestra qué conjugación respeta la taxonomía concreta de las
rutas y por qué partícula y antipartícula comparten banda.

\section{Persistencia finita y partícula virtual}

Sea
\[
\cdots\xrightarrow{\rho_{M+1,M}}S_M
\xrightarrow{\rho_{M,M-1}}S_{M-1}\xrightarrow{}\cdots
\]
el sistema de estados admisibles a profundidades finitas. Para
\(s_N\in S_N\), definimos el conjunto de prolongaciones
\[
\mathcal P_M(s_N)
:=
\{s_M\in S_M:\rho_{M,N}(s_M)=s_N\},
\qquad M\geq N.
\]
Su profundidad de vida es
\[
\mathcal L(s_N)
:=\sup\{M\geq N:\mathcal P_M(s_N)\neq\varnothing\}
\in\mathbb N\cup\{\infty\}.
\]

\begin{definition}[Persistencia y virtualidad]
\label{def:persistencia-virtualidad}
Un prefijo \(s_N\) es \emph{persistente} si pertenece a una sección global
\((s_M)_{M\geq N}\) compatible. Es \emph{virtual finito} si
\(\mathcal L(s_N)<\infty\), existe una prolongación hasta
\(\mathcal L(s_N)\) y no existe ninguna a la profundidad siguiente.
\end{definition}

\begin{theorem}[Criterio de persistencia]
\label{thm:criterio-persistencia-virtual}
Supóngase que cada estado posee un número finito de prolongaciones inmediatas
y que las restricciones son coherentes. Entonces:
\[
\mathcal L(s_N)=\infty
\quad\Longleftrightarrow\quad
s_N\text{ pertenece a una sección global}.
\]
Si \(\mathcal L(s_N)=L<\infty\), el prefijo transporta ruta, memoria,
registro y firma hasta \(L\), pero no define una partícula persistente.
\end{theorem}

\begin{proof}
Una sección global proporciona una prolongación a toda profundidad, de modo
que \(\mathcal L=\infty\). Recíprocamente, las prolongaciones de \(s_N\)
forman un árbol enraizado, infinito y finitamente ramificado. Por el lema de
Kőnig contiene una rama infinita, cuyos vértices son compatibles por
construcción y definen una sección global.

Si \(\mathcal L=L<\infty\), existe un estado a profundidad \(L\) y ninguno
a \(L+1\). Los lectores definidos sobre prefijos pueden evaluarlo hasta
\(L\); la ausencia de prolongación impide convertir esa firma finita en una
sección del límite inverso.
\end{proof}

La partícula virtual posee así una definición interna. No es una partícula
``menos real'' ni una violación breve de una ley de conservación. Es una
ruta cuya memoria influye durante una profundidad finita y cuya genealogía
se extingue antes de estabilizarse. El contraste con propagadores y estados
intermedios de la física convencional se realiza después de identificar qué
lector representa esas profundidades.

\section{Cadena rectora}

Los resultados de este capítulo fijan la dirección completa:
\[
\boxed{
\text{extensión}\to
\text{ruta}\to
\text{registro}\to
\text{firma}\to
\text{banda orientada}\to
\text{partícula}\to
\text{masa}.}
\]
La antipartícula es la segunda sección de la banda. La virtualidad es una
persistencia finita. Ninguna de estas nociones se define a partir de una
masa observada.

\def\HMTIncludeRouteLedgerBlock{1}
\input{sections/md/05b_extension_superviviente_caracter}
\let\HMTIncludeRouteLedgerBlock\undefined
\input{sections/md/21_persistencia_ruta_mobius_familias_actualizada}
\input{sections/md/22_bisagra_hojas_pi_phi2_actualizada}
\HMTRestoreHeadings

\chapter{Coordenadas angulares intrínsecas y canales orientados}
\label{chap:angulos-canales-rev7}
\HMTClaim{MD-R-ANGLE}
\HMTClaim{MD-R-CONTRAANGLE}
\HMTClaim{MD-R-QPLUS-QMINUS}
\HMTDemoteHeadings
\input{sections/md/02_pantalla_accion}
\input{sections/md/03_canales_contraangulo}
\HMTRestoreHeadings

\chapter{Transferencia a la recta de acción y cuantización por ciclo completo}
\label{chap:transferencia-accion-rev7}
\HMTClaim{HMT-R-ACTION-H-HBAR}
La rama determinantal ya fijó en la Parte~I la década interna
\(-34\). Aquí esa salida se transporta a la recta de acción y se combina
después con las rectas de duración y velocidad. Esta rama y la del carácter
adimensional son entradas hermanas de la realización de masa: ninguna es
salida de la otra y, en particular, \(\chi_{\rm full}\) no genera
\(\hbar_{\rm int}\).
\HMTDemoteHeadings
\input{sections/md/03b_escala_accion_determinantal}
\input{sections/md/03c_elipse_accion_holonomia_rev8}
\HMTRestoreHeadings

\chapter{Carácter positivo y jerarquía espectral infinita}
\label{chap:caracter-torres-rev7}
\HMTClaim{MD-R-CHARACTER}
\HMTClaim{MD-R-TOWER-SEMIGROUP}
\HMTClaim{MD-R-TOWER-RECURRENCE}
La firma de ruta \(\mathbb Z^5\) produce primero el carácter central. El
refinamiento global se extiende después sobre todo el semigrupo
\(\langle90,120\rangle\), sin identificar \(s_{270}\) con
\(\zeta_{270}=135\Delta_4^2\) ni borrar la procedencia distinta de
\(b_{120}\) y \(\eta_{120}\).
\HMTDemoteHeadings
\def\HMTIncludeCharacterBlock{1}
\input{sections/md/05b_extension_superviviente_caracter}
\let\HMTIncludeCharacterBlock\undefined
\input{sections/md/05c_factorizacion_split_caracter}
\input{sections/md/05a_jerarquia_espectral_infinita_rev7}
\input{sections/md/05_torres_globales}
\HMTRestoreHeadings

\chapter{Atlas matemático de rutas, familias y representaciones internas}
\label{chap:atlas-especies-rev7}
\HMTClaim{MD-R-ATLAS-81-56-13}
\HMTClaim{MD-R-COLOR}
El atlas es el clasificador estructural que recibe las rutas ya tipadas. No
sustituye la cadena ledger--firma--carácter ni se deduce de una tabla de
masas.
\HMTDemoteHeadings
\def\HMTIncludeAtlasStructural{1}
\input{sections/md/06h_atlas_familias_rev6}
\let\HMTIncludeAtlasStructural\undefined
\input{sections/md/06e_color_qutrit_gauge_finito}
\HMTRestoreHeadings

\chapter{Ley general de masas y realizaciones de especie}
\label{chap:ley-masas-rev7}
\HMTClaim{MD-R-MASS-LAW}
La ley se aplica al dominio ya construido. La ventana numérica de doce
especies es un testigo basal; el catálogo externo de masas y anchuras es un
contraste posterior. Ninguno de los dos define el atlas ni genera una ruta.
\HMTDemoteHeadings
\input{sections/md/05d_ley_masas_holografica_rev6}
\input{sections/md/05e_contraste_masas_rev7}
\HMTRestoreHeadings

\chapter{Realización física del atlas e inventario humano de sectores}
\label{chap:inventario-sectores-rev8}
\HMTClaim{MD-R-MASS-DOMAIN}
\HMTClaim{MD-R-FLAVOR}
La clasificación física se abre sólo después de construir persistencia,
atlas, firma, carácter y ley. Los veintitrés registros siguientes
son sectores expositivos tipados: no son veintitrés masas, familias o
multisecciones, y permanecen separados de los cardinales
\(81/56/13/324\) y de las dos colecciones de \(192\) clases.
\HMTDemoteHeadings
\def\HMTIncludeAtlasPhysical{1}
\input{sections/md/06h_atlas_familias_rev6}
\let\HMTIncludeAtlasPhysical\undefined
\HMTRestoreHeadings

\chapter{El electrón como cierre holográfico central}
\label{chap:electron-rev8}
\HMTClaim{MD-R-ELECTRON}
El electrón no se reduce a un valor basal. Su biografía reúne el centro
\((5,5)\), la ruta activa, la firma cruda, la firma par CT108, el origen
afín, el selector interno \((-22,+15)\), los seis factores del cierre, el
carácter y la torre; la carta MeV y CODATA aparecen únicamente al final.
\HMTDemoteHeadings
\input{sections/md/06f_biografia_electron_rev6}
\HMTRestoreHeadings

\chapter{Composición de especies y sectores topológicos}
\label{chap:composicion-topologica-rev7}
\HMTClaim{MD-R-COMPOSITION}
\HMTClaim{MD-R-TOPOLOGY}
\HMTClaim{MD-R-MAJORANA}
\HMTClaim{MD-R-FIBONACCI-ANYON}
\HMTDemoteHeadings
\input{sections/md/06g_composicion_especies_rev6}
\def\HMTIncludeCMBlock{1}
\def\HMTIncludeTopologicalBlock{1}
\ifdefined\HMTIncludeCMBlock
\chapter[Estructura compleja \texorpdfstring{\(\widetilde J_M^{\,2}=-I\)}{J M al cuadrado = -I} y sectores topológicos]
{Estructura compleja \texorpdfstring{\(\widetilde J_M^{\,2}=-I\)}{J M al cuadrado = -I} sobre el estrato fijo, involuciones dodecafásicas y sectores topológicos}
\label{chap:sectores-topologicos-no-omision}
\index{sector topológico}\index{Fibonacci, anillo basado}
\index{Majorana}\index{Ising}\index{grupo de trenzas}

La involución \(C_M\), el anillo basado asociado a \(F_{\rm av}\), los
caracteres de trenza y las categorías apuntadas pertenecen a estratos
algebraicos distintos. Se presentan con sus dominios y controles de
identificación para que una coincidencia nominal no sustituya al morfismo que
debería relacionarlos. El censo del atlas, la jerarquía espectral infinita
cuyas ocho alturas históricas son testigos iniciales, y las completaciones
estadísticas conservan las demostraciones de
\cref{thm:censo-atlas-81,thm:descenso-necesario-fav,%
thm:principio-exclusion-pauli,thm:dos-cuatro-pi-estadistica}.

Cada resultado declara directamente su dominio, su construcción y su
conclusión; ninguna coincidencia nominal modifica esas dependencias.

\section{Estrato periódico de la involución dodecafásica}
\label{sec:involucion-cm-censo}

El dominio \(\mathscr T_{12}^{\rm word}\), la conjugación
\(C_M\mathbf t=-\operatorname{rev}(\mathbf t)\), su carácter involutivo y el
censo \(\lvert\operatorname{Fix}(C_M)\rvert=3^6=729\) tienen su residencia
primaria en las
\cref{def:palabras-cm,thm:involucion-cm-censo}, dentro de la banda de
partícula. Aquí se añade únicamente la condición periódica necesaria para el
estrato topológico. Si \(r\) denota el desplazamiento cíclico
\((r\mathbf t)_m=t_{m-1}\), ponemos
\[
 \operatorname{Per}_6
 :=\{\mathbf t\in\mathscr T_{12}^{\rm word}:r^6\mathbf t=\mathbf t\}.
\]

\begin{proposition}[Censo del estrato periódico]
\label{prop:estrato-periodico-cm-rev7}
Al imponer periodicidad de período divisor de seis,
\[
 \bigl|\operatorname{Fix}(C_M)\cap\operatorname{Per}_6\bigr|
 =3^3=27.
\]
Esta identidad es exacta.
\end{proposition}

\begin{proof}
El \cref{thm:involucion-cm-censo} da la forma general de los puntos fijos.
Si además \(r^6\mathbf t=\mathbf t\), la segunda mitad de la palabra coincide
con la primera. Escribiendo las seis primeras coordenadas, las dos
condiciones equivalen a
\[
 (t_0,t_1,t_2,t_3,t_4,t_5)
 =(a,b,c,-c,-b,-a).
\]
Los parámetros \(a,b,c\) son independientes y toman tres valores; por tanto,
la intersección tiene \(3^3\) elementos.
\end{proof}

El ejemplo general de un punto fijo, antes de exigir período seis, es
\[
 (a,b,c,d,e,f,-f,-e,-d,-c,-b,-a).
\]
Esta forma coordinada permite comprobar a mano tanto la involución como el
conteo y evita interpretar el entero \(729\) mediante otro objeto homónimo.

\begin{remark}[Censo exhaustivo de la involución]
\label{rem:procedencia-n40-cm}
La enumeración de las \(3^{12}=531\,441\) palabras no encuentra ninguna
violación de la involución y obtiene
\[
 |\operatorname{Fix}(C_M)|=729,
 \qquad
 |\operatorname{Fix}(C_M)\cap\operatorname{Per}_6|=27.
\]
La demostración primaria y la \cref{prop:estrato-periodico-cm-rev7} explican
ambos números sin depender del recorrido: los puntos fijos tienen seis
coordenadas libres, y la condición adicional de período seis reduce la
elección a tres. El censo exhaustivo confirma que no
existen órbitas excepcionales fuera de esa descripción.
\end{remark}

\begin{criterion}[Falsadores del censo \(C_M\)]
\label{crit:falsadores-cm}
Refutaría el resultado cualquiera de los hechos siguientes: una palabra para
la cual \(C_M^2\mathbf t\ne\mathbf t\); un punto fijo que no posea la forma
coordenada anterior; un total exhaustivo distinto de \(729\); o un total
distinto de \(27\) después de imponer \(r^6\mathbf t=\mathbf t\). La
coincidencia nominal de este \(729\) con otro cardinal no identifica sus
dominios.
\end{criterion}

\section{Estructura cuadrática transportada al estrato fijo}
\label{sec:estructura-cuadratica-estrato-fijo-cm}

Identificamos el alfabeto de signos con el cuerpo ternario mediante
\[
 \{-1,0,+1\}\xrightarrow{\sim}\mathbb F_3,
 \qquad
 -1\longmapsto2,\quad 0\longmapsto0,\quad +1\longmapsto1.
\]
Con esta identificación,
\(\mathscr T_{12}^{\rm word}\cong\mathbb F_3^{12}\), y tanto
\(C_M=-\operatorname{rev}\) como el desplazamiento \(r\) son operadores
\(\mathbb F_3\)-lineales. En consecuencia,
\(\operatorname{Fix}(C_M)\) y
\(\operatorname{Per}_6=\ker(r^6-I)\) son subespacios vectoriales.

La forma coordinada de los puntos fijos define el isomorfismo lineal
\[
 s_{\rm fix}:\mathbb F_3^6\xrightarrow{\sim}\operatorname{Fix}(C_M),
\]
\[
 s_{\rm fix}(a,b,c,d,e,f)
 =(a,b,c,d,e,f,-f,-e,-d,-c,-b,-a).
\]
Sea \(J_W\) el operador de Paley--Witt construido en el
\cref{thm:f9-tres-desde-paley}, con la base, la polarización y la orientación
allí declaradas.

\begin{proposition}[Estructura compleja \(\widetilde J_M^{\,2}=-I\) sobre el estrato fijo]
\label{prop:raiz-menos-uno-estrato-fijo-cm}
El transporte
\[
\widetilde J_M:=s_{\rm fix}J_Ws_{\rm fix}^{-1}
\in\operatorname{End}_{\mathbb F_3}
   \bigl(\operatorname{Fix}(C_M)\bigr)
\]
satisface
\[
 \widetilde J_M^{\,2}=-I.
\]
\end{proposition}

\begin{proof}
Por conjugación,
\[
 \widetilde J_M^{\,2}
 =s_{\rm fix}J_Ws_{\rm fix}^{-1}
  s_{\rm fix}J_Ws_{\rm fix}^{-1}
 =s_{\rm fix}J_W^2s_{\rm fix}^{-1}
 =-I.
\]
\end{proof}

El resultado está demostrado con la carta \(s_{\rm fix}\), la base, la polarización y
la orientación como estructura de partida explícita. No identifica
\(C_M\) con \(J_W\): el primero es una involución del espacio de palabras y
el segundo una estructura cuadrática sobre su estrato fijo. Tampoco demuestra
que el transporte sea natural respecto de toda la dinámica TPK.

Pongamos
\[
 W_{27}
 :=\operatorname{Fix}(C_M)\cap\operatorname{Per}_6.
\]
La inclusión
\[
 W_{27}\hookrightarrow\operatorname{Fix}(C_M)
\]
es canónica, y la \cref{prop:estrato-periodico-cm-rev7} prueba que
\(\dim_{\mathbb F_3}W_{27}=3\).

\begin{corollary}[Obstrucción cuadrática en el estrato de veintisiete]
\label{cor:obstruccion-cuadratica-w27}
El subespacio \(W_{27}\) no puede ser estable bajo
\(\widetilde J_M\).
\end{corollary}

\begin{proof}
Si fuera estable, la relación
\(\widetilde J_M^2=-I\) lo convertiría en un módulo sobre
\[
 \mathbb F_3[X]/(X^2+1)\cong\mathbb F_9.
\]
Todo espacio vectorial sobre \(\mathbb F_9\) tiene dimensión par cuando se
considera sobre \(\mathbb F_3\), en contradicción con
\(\dim_{\mathbb F_3}W_{27}=3\).
\end{proof}

La igualdad cardinal entre \(W_{27}\) y las veintisiete rectas de Schläfli
no proporciona un mapa entre ambos objetos. Tal relación exigiría construir
una aplicación de incidencia que preserve la orientación central y la acción
del estabilizador, además de una clausura dinámica compatible. No se
interpone el atlas de \(81\) celdas entre \(W_{27}\) y el estrato fijo de
\(729\) estados: el atlas de partículas pertenece a otro dominio.

\section[Niveles de Majorana]{Estratificación matemática del objeto Majorana en cuatro niveles}
\label{sec:cuatro-niveles-majorana}

La inversión \(C_M\), la paridad exterior, un fermión relativista de
Majorana y un modo cero de Majorana son objetos de tipos distintos:

\begin{enumerate}
 \item \(C_M=-\operatorname{rev}\) actúa sobre
 \(\mathscr T_{12}^{\rm word}\). Su involutividad y su censo se demuestran
 exactamente; esta construcción no proporciona por sí misma un espinor, una
 ecuación de campo ni un término de masa.

 \item El carácter de signo \(q=-1\) y la completación exterior pertenecen
 al álgebra de intercambio. Relativamente a CAR y a la graduación declarada,
 producen \(N_s^2=N_s\), espectro \(\{0,1\}\) y el retorno de orientación
 \(2\pi/4\pi\). Estas identidades se deducen de CAR y de la graduación; no se
 deducen del
 mero censo de \(C_M\).

 \item Un fermión relativista de Majorana es un objeto espinorial provisto de
 representación relativista, conjugación de carga compatible, condición de
 autorconjugación, término de masa y dinámica. Esos datos no forman parte de
 la definición de \(C_M\) ni se deducen de su censo.

 \item Un modo cero de Majorana, o equivalentemente el sector no abeliano de
 una teoría de Ising en la acepción topológica, requiere simples
 \(\mathbf 1,\psi,\sigma\) con
 \[
  \psi\otimes\psi\simeq\mathbf 1,\qquad
  \psi\otimes\sigma\simeq\sigma,\qquad
  \sigma\otimes\sigma\simeq\mathbf 1\oplus\psi,\qquad
  d_\sigma=\sqrt2,
 \]
 junto con asociador y trenzado coherentes. La brecha, el desdoblamiento y la
 plataforma pertenecen a una realización física adicional; las reglas de
 fusión anteriores no definen una masa de polo universal.
\end{enumerate}

La separación es de dominio y no de nomenclatura: involución de palabras,
paridad CAR, espinor autorconjugado y simple no abeliano son cuatro objetos
matemáticos distintos. En ausencia de un morfismo explícito entre ellos, no
se identifica ninguno por compartir el nombre «Majorana».

\fi

\ifdefined\HMTIncludeTopologicalBlock
\chapter[Incidencia y sectores topológicos]{Incidencia, Fibonacci y sectores topológicos}
\label{chap:incidencia-fibonacci-topologia}
\index{sector topológico}\index{Fibonacci, anillo basado}
\index{grupo de trenzas}

La incidencia residuo--soporte prolonga el atlas hacia estructuras de fusión
y trenza. El orden es esencial: primero se construye el anillo basado
Fibonacci y sus dimensiones; después se estudian las representaciones
algebraicas de trenzas y el centro apuntado. Cada objeto conserva su propio
dominio y no se identifica con una partícula física por compartir nombre o
cardinal.

\section{Incidencia residuo--soporte}
\label{sec:incidencia-residuo-soporte}

Fijamos las clases de residuo y los soportes en el orden
\[
 \begin{aligned}
 C_1&=\{1,4,7\},&
 C_2&=\{2,5,8\},&
 C_0&=\{3,6,9\},\\
 \mathcal D_{\rm act}&=\{1,3,5,7\},&
 \mathcal D_{\rm evt}&=\{2,4,6,8\},&
 \mathcal D_{\bar0}^{(9)}&=\{9\}.
 \end{aligned}
\]
El último símbolo significa «soporte del residuo nonádico nulo». No designa
el dominio completo \(D_9=\{1,\ldots,9\}\) ni un sector físico de masa nula.

Sean \(E_{\mathcal D}\cong\mathbb Z^3\) y
\(E_C\cong\mathbb Z^3\) los módulos libres con bases ordenadas
\[
 (\mathcal D_{\rm act},\mathcal D_{\rm evt},
   \mathcal D_{\bar0}^{(9)})
 \quad\text{y}\quad
 (C_1,C_2,C_0),
\]
respectivamente. La aplicación de incidencia
\(B_{\rm res,sup}:E_{\mathcal D}\to E_C\) cuenta intersecciones.

\begin{theorem}[Forma de Smith e imagen de la incidencia]
\label{thm:incidencia-residuo-soporte}
En las bases precedentes,
\[
 B_{\rm res,sup}
 =\bigl(|C_i\cap\mathcal D_j|\bigr)_{i,j}
 =
 \begin{pmatrix}
 2&1&0\\
 1&2&0\\
 1&1&1
 \end{pmatrix}.
\]
Además,
\[
 \det B_{\rm res,sup}=3,\qquad
 \operatorname{SNF}(B_{\rm res,sup})
 =\operatorname{diag}(1,1,3),
\]
\[
 \operatorname{im}B_{\rm res,sup}
 =\{(u,v,w)\in\mathbb Z^3:u+v\equiv0\pmod3\}.
\]
Si
\[
 P=\begin{pmatrix}0&1&0\\1&0&0\\0&0&1\end{pmatrix},
\]
el intercambio simultáneo
\(C_1\leftrightarrow C_2\) y
\(\mathcal D_{\rm act}\leftrightarrow\mathcal D_{\rm evt}\) es equivariante:
\[
 P B_{\rm res,sup}P=B_{\rm res,sup}
 \qquad\Longleftrightarrow\qquad
 PB_{\rm res,sup}=B_{\rm res,sup}P.
\]
\end{theorem}

\begin{proof}
Las nueve entradas proceden de las intersecciones
\[
\begin{array}{cccc}
 &\mathcal D_{\rm act}&\mathcal D_{\rm evt}
 &\mathcal D_{\bar0}^{(9)}\\
 C_1&\{1,7\}&\{4\}&\varnothing\\
 C_2&\{5\}&\{2,8\}&\varnothing\\
 C_0&\{3\}&\{6\}&\{9\}.
\end{array}
\]
La expansión por la tercera columna da
\[
 \det B_{\rm res,sup}
 =\det\begin{pmatrix}2&1\\1&2\end{pmatrix}=3.
\]
El máximo común divisor de las entradas es uno. El menor de orden dos
formado por las filas primera y tercera y las columnas segunda y tercera
tiene determinante uno:
\[
 \det\begin{pmatrix}1&0\\1&1\end{pmatrix}=1.
\]
Por los divisores de Smith, \(d_1=d_2=1\). El producto
\(d_1d_2d_3\) coincide con el valor absoluto del determinante, que el
cálculo anterior fija en tres. Por consiguiente,

\[
 d_1d_2d_3=3.
\]
La forma normal es, por tanto, \(\operatorname{diag}(1,1,3)\).

Cada columna satisface \(u+v\equiv0\pmod3\), así que la imagen está contenida
en el submódulo indicado. Recíprocamente, si \(u+v\equiv0\pmod3\), entonces
\[
 x=\frac{2u-v}{3},\qquad
 y=\frac{2v-u}{3},\qquad
 z=w-x-y
\]
son enteros y
\[
 B_{\rm res,sup}(x,y,z)^{\mathsf T}=(u,v,w)^{\mathsf T}.
\]
Esto prueba la igualdad de imágenes. Finalmente, la multiplicación explícita
por \(P\) intercambia las dos primeras filas y las dos primeras columnas y
deja la matriz invariante.
\end{proof}

\begin{criterion}[Control de la incidencia]
\label{crit:incidencia-residuo-soporte}
Falsaría el teorema una sola intersección mal contada, un menor que alterase
los divisores de Smith, un vector de la imagen con \(u+v\not\equiv0\pmod3\),
un vector que satisficiera la congruencia pero no admitiese la preimagen
exhibida, o la pérdida de la equivarianza \(PBP=B\).
\end{criterion}

\section[Anillo basado de Fibonacci]{Anillo basado asociado a la incidencia
  \texorpdfstring{\(F_{\rm av}\)}{F\_av}}
\label{sec:anillo-basado-fibonacci}

El \cref{thm:descenso-necesario-fav} demuestra, a partir de la órbita de
modos \((\Sigma,\Pi,\Pi)\), la matriz
\[
 F_{\rm av}=
 \begin{pmatrix}0&1\\1&1\end{pmatrix}.
\]
De esta matriz se obtiene el siguiente anillo basado.

\begin{definition}[Anillo basado de Fibonacci]
\label{def:anillo-basado-fibonacci}
Sea
\[
 K_{\rm Fib}:=\mathbb Z\mathbf 1\oplus\mathbb Z\tau
\]
el grupo abeliano libre con unidad \(\mathbf 1\) y producto bilineal fijado por
\[
 \tau\cdot\mathbf 1=\tau,\qquad
 \tau\cdot\tau=\mathbf 1+\tau.
\]
En este estrato, \(+\) y \(\cdot\) son la adición y la multiplicación del
anillo. Los símbolos categóricos \(\oplus\) y \(\otimes\) se reservan para
una categorificación elegida posteriormente.
\end{definition}

\begin{theorem}[Realización de \(F_{\rm av}\) y dimensión de Perron]
\label{thm:anillo-fibonacci-fav}
En la base ordenada \((\mathbf 1,\tau)\), la multiplicación por \(\tau\) en
\(K_{\rm Fib}\) tiene matriz \(F_{\rm av}\). Si
\(F_0=0\), \(F_1=1\) y
\(F_{n+1}=F_n+F_{n-1}\) para \(n\geq1\), entonces
\[
 \tau^n=F_{n-1}\mathbf 1+F_n\tau
 \qquad(n\geq1).
\]
El único homomorfismo positivo de dimensión que envía
\(\mathbf 1\) a uno satisface
\[
 \operatorname{FPdim}(\tau)=\varphi.
\]
La identificación de las dos hojas tipadas con la base
\((\mathbf 1,\tau)\) determina por completo este anillo y su dimensión
positiva.
\end{theorem}

\begin{proof}
Las columnas de la matriz de multiplicación por \(\tau\) son las coordenadas
de
\[
 \tau\cdot\mathbf 1=\tau
 \quad\text{y}\quad
 \tau\cdot\tau=\mathbf 1+\tau;
\]
por tanto son \((0,1)^{\mathsf T}\) y \((1,1)^{\mathsf T}\), respectivamente.
Se obtiene así \(F_{\rm av}\).

La fórmula de potencias vale para \(n=1\). Si vale para \(n\), entonces
\[
 \tau^{n+1}
 =F_{n-1}\tau+F_n(\mathbf 1+\tau)
 =F_n\mathbf 1+F_{n+1}\tau,
\]
lo que cierra la inducción. Finalmente, una dimensión multiplicativa positiva
\(d=\operatorname{FPdim}(\tau)\) obedece
\[
 d^2=1+d.
\]
Sus raíces son \(\varphi\) y \(-\varphi^{-1}\); la positividad selecciona
\(d=\varphi\).
\end{proof}

La denominación \emph{Fibonacci} queda así explicada por la propia regla de
fusión: al iterar \(\tau\), los dos canales posibles aparecen con
multiplicidades consecutivas \(F_{n-1}\) y \(F_n\). El espacio de estados de
fusión crece, por tanto, con razón asintótica \(\varphi\). Ésta es la base
combinatoria de su uso computacional: la información se codifica en canales
globales de fusión y no en una masa asignada a un punto.

\begin{criterion}[Condiciones de una categorización Fibonacci]
\label{crit:categorificacion-fibonacci}
La capa anular quedaría refutada si la matriz de multiplicación no fuese
\(F_{\rm av}\), si fallase la recurrencia de potencias o si la dimensión
positiva no fuese \(\varphi\). Una categorización trenzada cuyo anillo de
Grothendieck sea el de la \cref{def:anillo-basado-fibonacci} exige además un
funtor monoidal
\[
 \mathfrak C_{\rm Fib}:\mathsf{TPK}
 \dashrightarrow\mathcal C_{\rm Fib},
\]
la verificación de pentágono y hexágono en su imagen y una regla interna que
seleccione calibre y quiralidad.
Estas condiciones no son consecuencias de la identidad anular por sí sola.
\end{criterion}

\section[Caracteres de trenza]{Seis caracteres y cuatro torsiones de trenza}
\label{sec:rescate-trenzas-ex02}

Sea \(B_n\), \(n\geq2\), el grupo de trenzas de Artin con generadores
\(\sigma_1,\ldots,\sigma_{n-1}\), relaciones
\[
 \sigma_i\sigma_{i+1}\sigma_i
 =\sigma_{i+1}\sigma_i\sigma_{i+1},
 \qquad
 \sigma_i\sigma_j=\sigma_j\sigma_i\quad(|i-j|\geq2),
\]
y sea \(w:B_n\to\mathbb Z\) la suma de exponentes. Los lectores
\[
 w_2:=w\bmod2,\qquad w_3:=w\bmod3
\]
se reúnen, por el teorema chino del resto, en
\[
 w_6:B_n\longrightarrow
 \mathbb Z/2\mathbb Z\times\mathbb Z/3\mathbb Z
 \simeq\mathbb Z/6\mathbb Z.
\]

\begin{theorem}[Abelianización, seis caracteres y torsiones de permutación]
\label{thm:trenzas-caracteres-z6}
Para \(n\geq2\),
\[
 B_n^{\rm ab}\simeq\mathbb Z,
\]
mediante \(w\). En consecuencia, \(w_6\) produce seis caracteres unitarios
\[
 \chi_k(b)
 :=\exp\!\left(\frac{2\pi\mathrm i k}{6}w(b)\right),
 \qquad k\in\mathbb Z/6\mathbb Z.
\]

Sea \(V\ne0\) un espacio de Hilbert, sea
\(P_i:V^{\otimes n}\to V^{\otimes n}\) el intercambio de los factores
\(i,i+1\), y sea \(q\in\mu_6\). La regla
\[
 \rho_q(\sigma_i):=qP_i,
 \qquad
 \rho_q(b)=q^{w(b)}\pi_n(b),
\]
donde \(\pi_n\) es la representación de permutación, define una
representación de \(B_n\). Esta representación desciende al grupo simétrico
si y sólo si
\[
 q^2=1.
\]
Para los cuatro valores \(q\in\mu_6\setminus\{\pm1\}\),
\[
 \operatorname{Inv}(\rho_q)=0.
\]
Esos cuatro valores definen torsiones escalares genuinamente trenzadas, pero
no constituyen por sí solos cuatro espacios de Fock ni cuatro sectores
anyónicos físicos.
\end{theorem}

\begin{proof}
En todo cociente abeliano, la relación de trenza da
\[
 2[\sigma_i]+[\sigma_{i+1}]
 =[\sigma_i]+2[\sigma_{i+1}],
\]
luego todos los generadores tienen la misma clase. No queda otra relación
sobre esa clase, pues \(w(\sigma_i)=1\); por tanto la abelianización es
\(\mathbb Z\).

Los operadores \(P_i\) satisfacen las relaciones de trenza y
\(P_i^2=I\). Multiplicar cada generador por el mismo escalar \(q\) conserva
las dos relaciones, de modo que \(\rho_q\) es una representación. Para
factorizar a través de \(S_n\) debe satisfacer la relación adicional
\(\sigma_i^2=1\), y
\[
 \rho_q(\sigma_i)^2=q^2I.
\]
Esto prueba el criterio \(q^2=1\).

Si \(v\) fuera invariante, entonces
\(qP_i v=v\), es decir, \(P_i v=q^{-1}v\), para todo \(i\). Como cada
involución \(P_i\) sólo tiene autovalores \(\pm1\), no existe tal vector no
nulo cuando \(q\ne\pm1\). Por consiguiente,
\(\operatorname{Inv}(\rho_q)=0\) para los cuatro valores restantes.
\end{proof}

Cuando \(\dim V>1\) y \(n\geq3\), los operadores \(P_1\) y \(P_2\) no
conmutan; por ello la representación completa no se vuelve abeliana por
proceder de un carácter escalar. El carácter \(\chi_k\) y la representación
\(\rho_q\) son objetos distintos: el primero es unidimensional; la segunda
retiene la acción de permutación.

\section[Centro modular y realizaciones topológicas]{Centro modular apuntado y condiciones de realización Fibonacci e Ising}
\label{sec:centro-apuntado-no-go}

Sea
\[
 A=(\mathbb Z/9\mathbb Z)^2
\]
y sea \(\operatorname{Vec}_A\) la categoría de espacios vectoriales
\(A\)-graduados con asociador trivial. Fijar un bicharácter no degenerado
identifica \(\widehat A\) con \(A\), pero no altera los conteos siguientes.

\begin{theorem}[Centro apuntado y obstrucción dimensional]
\label{thm:centro-apuntado-no-go}
El centro de Drinfeld
\[
 \mathcal Z(\operatorname{Vec}_A)
\]
es una categoría modular apuntada cuyos simples se parametrizan por
\((a,\chi)\in A\times\widehat A\). Por tanto posee
\[
 |A|\,|\widehat A|=81^2=6561
\]
simples, todos de dimensión cuántica uno, y dimensión cuántica total
\[
 \mathcal D
 =\sqrt{\sum_{x\ {\rm simple}}d_x^2}
 =\sqrt{6561}=81.
\]
En esta categoría no puede realizarse, preservando simples y dimensiones de
Frobenius--Perron, ni el anillo Fibonacci ni el sector Ising: éstos requieren
respectivamente un simple de dimensión \(\varphi\) y un simple de dimensión
\(\sqrt2\).
\end{theorem}

\begin{proof}
Para un grupo abeliano finito y asociador trivial, cada objeto simple del
centro queda determinado por un grado \(a\in A\) y un carácter
\(\chi\in\widehat A\) que especifica su media trenza. Tanto el objeto
graduado como el carácter son unidimensionales; en consecuencia, todos los
simples son invertibles y tienen dimensión cuántica uno. Como
\(|A|=|\widehat A|=9^2=81\), el número de pares es \(81^2\), y la fórmula de
\(\mathcal D\) da \(81\).

Toda subcategoría de fusión de una categoría apuntada vuelve a ser apuntada;
en particular, sus simples tienen dimensión uno. Esto contradice
\(\operatorname{FPdim}(\tau)=\varphi\) para Fibonacci y
\(d_\sigma=\sqrt2\) para Ising. La obstrucción es dimensional y no depende
de una semejanza de nombres o cardinales.
\end{proof}

\begin{criterion}[Condiciones para una realización física de los sectores]
\label{crit:promocion-sectores-topologicos}
Los resultados de
\cref{thm:trenzas-caracteres-z6,thm:centro-apuntado-no-go} quedarían
algebraicamente falsados por un fallo de las relaciones de Artin, del
criterio \(q^2=1\), del censo \(6561\), de la apuntabilidad o de
\(\mathcal D=81\). Una realización física exige, además, un
Hamiltoniano o protocolo impulsado, una brecha espectral, operadores de
trenza realizables, protección frente a perturbaciones y un mapa desde las
multisecciones enriquecidas. Ninguno de esos datos se deduce únicamente de
los dos teoremas algebraicos anteriores.
\end{criterion}
\fi

\let\HMTIncludeTopologicalBlock\undefined
\let\HMTIncludeCMBlock\undefined
\HMTRestoreHeadings

\chapter{Representaciones cuánticas de rutas}
\label{chap:estadistica-cuantica-rev7}
\HMTClaim{MD-R-STATISTICS}
\HMTClaim{MD-R-NULL-SECTOR}
\HMTClaim{MD-R-ANTIPARTICLE}
\HMTDemoteHeadings
\def\HMTIncludeParticleQuantum{1}
\input{sections/md/04_particula_estadistica}
\let\HMTIncludeParticleQuantum\undefined
\input{sections/md/04d_ocupacion_estadistica}
\input{sections/md/delta_127/04b_completaciones_cuanticas}
% Integración por delta de la adenda de realizaciones físicas.
% Residencia principal: transferencia física de las tres álgebras y lector temporal.

\chapter{Regímenes cuadráticos y lectura temporal}
\label{chap:regimenes-cuadraticos-tiempo}
\index{álgebra cuadrática!tres regímenes}
\index{cristal de tiempo!realización finita HMT}
\index{tiempo!lector TRIT}

\section{Álgebras cuadráticas del TRIT y signaturas de curvatura}
\label{sec:regimenes-cuadraticos-recibidos}

La construcción y la demostración primarias de los tres regímenes residen en
el \cref{chap:trit-algebras}. Para
\(\tau\in\TRIT\), con imagen balanceada
\(\bar\tau\in\{-1,0,+1\}\), se tiene
\[
 \mathbb A_\tau=\RR[X]/(X^2+\bar\tau),
\]
y, si \(J\) es la clase de \(X\),
\[
\begin{array}{ccl}
 \tau=+1&:&J^2=-I,\quad\text{régimen elíptico},\\
 \tau=0&:&J^2=0,\quad\text{régimen parabólico},\\
 \tau=-1&:&J^2=+I,\quad\text{régimen hiperbólico}.
\end{array}
\]
Para los tres valores de \(\tau\), el cociente
\(\mathbb A_\tau=\mathbb R[X]/(X^2+\bar\tau)\) satisface exactamente las tres
identidades cuadráticas anteriores. En particular, la raíz interna de
\(-1\) se realiza como operador en la hoja elíptica; HMT no necesita negar ni
suprimir la extensión compleja. El álgebra parabólica contiene un nilpotente,
y la hiperbólica se descompone mediante los dos idempotentes
\((I\pm J)/2\).

Estas identidades clasifican álgebras y transportes. No seleccionan por sí
solas una métrica espaciotemporal ni prueban que el universo físico ocupe una
de las tres hojas.

\section{Familia métrica parametrizada del lector temporal}
\label{sec:ansatz-metrico-tres-curvaturas}

Escribamos
\[
 \mathbb A_s:=\RR[X]/(X^2+s),
 \qquad s\in\{+1,0,-1\},
\]
para la misma familia con el signo real como índice.
Para hacer precisa la interfaz posible, fijemos una escala
\(\rho_{\rm curv}>0\), de dimensión \(L^{-1}\), y
\(\theta\in\RR/2\pi\ZZ\). Para \(s\in\{+1,0,-1\}\), sea
\[
 g_{s,\rho}
 =\dd r^2+S_{s,\rho}(r)^2\,\dd\theta^2,
\]
\[
 S_{s,\rho}(r)=
 \begin{cases}
  \rho_{\rm curv}^{-1}\sin(\rho_{\rm curv}r),&s=+1,\\
  r,&s=0,\\
  \rho_{\rm curv}^{-1}\sinh(\rho_{\rm curv}r),&s=-1.
 \end{cases}
\]
El dominio radial es \(r\geq0\) para \(s=0,-1\); para \(s=+1\), la carta
polar cubre \(0\leq r\leq\pi/\rho_{\rm curv}\), con la regularidad polar
usual en los extremos. En el interior de cada carta,
\[
 K_{\rm G}
 =-\frac{S_{s,\rho}''}{S_{s,\rho}}
 =s\,\rho_{\rm curv}^{\,2}.
\]
La última igualdad es un cálculo exacto del modelo una vez fijada la
realización. La flecha
\[
 \mathbb A_s\dashrightarrow(M_s,g_{s,\rho})
\]
es, sin embargo, un \emph{ansatz} de \textsc{interfaz física}; no es una
consecuencia canónica de la relación \(J^2=-sI\).

Si \(\ell_0\in\mathcal L_L^+\) es una longitud y se declara
\[
 \rho_{\rm curv}:=\frac{\Delta_4}{\ell_0},
\]
entonces
\[
 K_{\rm G}
 =s\,\frac{\Delta_4^2}{\ell_0^2}.
\]
Esta especialización está
\textsc{demostrada con estructura de partida explícita}: la conclusión se
obtiene una vez declarados \(s,\Delta_4,\ell_0\), pero no selecciona
\(\ell_0\), no
deduce la curvatura cosmológica observada ni identifica las razones
\(\mathcal G_{\rm tor}/\mathcal I_{\rm tor}\) o
\(\mathcal G_{\rm av}/\mathcal I_{\rm av}\) con \(K_{\rm G}\).

\section{Lector temporal del TRIT sobre ({-1,0,+1}): apertura, umbral y retorno con memoria}
\label{sec:lector-temporal-trit}

Una vez elegida una parametrización ordenada de las transiciones, el lector
temporal activo es
\[
\begin{aligned}
 +1&\longmapsto
 \text{retorno, memoria retenida y pasado},\\
 0&\longmapsto
 \text{umbral de evaluación y presente},\\
 -1&\longmapsto
 \text{apertura bidireccional y futuro}.
\end{aligned}
\]
El lector está \textsc{demostrado con estructura de partida explícita}: sus
conclusiones se obtienen una vez declarados la firma TRIT, el orden de la
trayectoria y la memoria transportada. No convierte el conjunto
\(\{-1,0,+1\}\) en una ecuación de movimiento. El futuro designa el espacio
de prolongaciones compatibles; la supervivencia puede seleccionar una rama
o una multisección sin borrar la genealogía de las restantes.

La frecuencia
\[
 \omega_0=\frac{2\pi}{108\,t_0}
\]
tipa el reloj del retorno completo mediante la entropía de decisión
\(I_{\mathrm{TRIT}}=\log 3\) y el transductor de Boltzmann.
No demuestra, aislada, una fase subarmónica estable. Sí permite construir el
Hamiltoniano finito de referencia que sigue; la selección de un Hamiltoniano
físico robusto requiere condiciones adicionales.

\section{Dinámica unitaria en 108 estados: Hamiltoniano, propagador y período}
\label{sec:hamiltoniano-reloj-108}

Sea
\[
 \mathcal H_{\rm clk}:=\mathbb C^{108},
 \qquad
 \{|j\rangle:j\in\mathbb Z/108\mathbb Z\}
\]
con su producto hermítico canónico. Fijemos
\(\zeta:=\exp(2\pi\mathrm i/108)\), el desplazamiento unitario
\[
 U_{\rm clk}|j\rangle=|j+1\rangle
\]
y la base de Fourier
\[
 |\widetilde k\rangle
 :=\frac1{\sqrt{108}}\sum_{j=0}^{107}\zeta^{jk}|j\rangle,
 \qquad 0\leq k\leq107.
\]
Con \(t_0\in\mathcal L_T^+\) y
\(\hbar_{\rm int}\in\mathcal L_S^+\), definimos
\[
 \boxed{
 H_{\rm clk}
 :=\hbar_{\rm int}\omega_0
 \sum_{k=0}^{107}k\,
 |\widetilde k\rangle\langle\widetilde k|,
 \qquad
 \omega_0=\frac{2\pi}{108\,t_0}.}
 \tag{CLK1}\label{eq:hamiltoniano-reloj-108}
\]
El operador tiene tipo energético
\([H_{\rm clk}]=\mathcal L_S\mathcal L_T^{-1}=\mathcal L_E\).

\begin{theorem}[Realización unitaria exacta del calendario \(108\)]
\label{thm:reloj-unitario-108}
El operador \(H_{\rm clk}\) es autoadjunto y su propagador durante un paso
satisface
\[
 U_{\rm step}
 :=\exp\!\left(-\frac{\mathrm i}{\hbar_{\rm int}}
 H_{\rm clk}t_0\right)
 =U_{\rm clk},
 \qquad
 U_{\rm step}^{108}=I.
 \tag{CLK2}\label{eq:propagador-reloj-108}
\]
Si
\[
 Z_{\rm clk}|j\rangle=\zeta^j|j\rangle,
\]
entonces, para cualquier estado de base \(|j_0\rangle\),
\[
 \langle Z_{\rm clk}\rangle_{n}
 :=\langle j_0|U_{\rm step}^{-n}Z_{\rm clk}
 U_{\rm step}^n|j_0\rangle
 =\zeta^{j_0+n}.
 \tag{CLK3}\label{eq:respuesta-reloj-108}
\]
La sucesión posee periodo primitivo \(108\).
\end{theorem}

\begin{proof}
La expresión \eqref{eq:hamiltoniano-reloj-108} es una descomposición
espectral con valores reales, luego \(H_{\rm clk}=H_{\rm clk}^\dagger\).
Además,
\[
 U_{\rm clk}|\widetilde k\rangle
 =\zeta^{-k}|\widetilde k\rangle,
\]
mientras que la exponencial de
\(\hbar_{\rm int}\omega_0 k\) durante \(t_0\) produce
\(\exp(-2\pi\mathrm ik/108)=\zeta^{-k}\). Ambos operadores coinciden sobre
una base ortonormal. Las identidades restantes siguen de
\(U_{\rm step}^n|j_0\rangle=|j_0+n\rangle\) y de la primitividad de
\(\zeta\).
\end{proof}

La formulación variacional correspondiente, para una curva
\(\psi\in C^1([t_a,t_b],\mathcal H_{\rm clk})\), no restringida durante la
variación, es
\[
 \boxed{
 L_{\rm clk}(\psi,\dot\psi)
 =\frac{\mathrm i\hbar_{\rm int}}2
 \bigl(
 \langle\psi|\dot\psi\rangle
 -\langle\dot\psi|\psi\rangle
 \bigr)
 -\langle\psi|H_{\rm clk}|\psi\rangle.}
 \tag{CLK4}\label{eq:lagrangiano-reloj-108}
\]
El funcional de acción es
\[
 \boxed{
 S_{\rm clk}[\psi]
 :=\int_{t_a}^{t_b}
 L_{\rm clk}(\psi(t),\dot\psi(t))\,\mathrm dt,}
 \tag{CLK5}\label{eq:accion-reloj-108}
\]
con \(\delta\psi(t_a)=\delta\psi(t_b)=0\).
\index{Schrödinger!ecuación del reloj finito}
Variar \(\psi\) y \(\bar\psi\) independientemente produce la ecuación de
Schrödinger del reloj finito
\[
 \boxed{\mathrm i\hbar_{\rm int}\dot\psi=H_{\rm clk}\psi}
 \tag{CLK6}\label{eq:schrodinger-reloj-108}
\]
junto con su ecuación adjunta \cite{SchrodingerNobel}. Su propagador muestreado
en \(t=nt_0\) es \(\psi_n=U_{\rm step}^{n}\psi_0\). Como \(H_{\rm clk}\) es
autoadjunto, la norma se conserva; un dato inicial normalizado permanece
normalizado. El resultado es exacto respecto de
\(H_{\rm clk},t_0,\hbar_{\rm int}\) declarados y no selecciona un Hamiltoniano
físico ni introduce una frecuencia experimental.

El \cref{thm:reloj-unitario-108} demuestra la dinámica unitaria del calendario
respecto de la completación compleja, \(t_0\) y \(\hbar_{\rm int}\)
declarados. Las cinco condiciones dinámicas de la sección siguiente reciben
una realización constructiva en \cref{sec:hmtf2-cristal-temporal}: allí se
fijan el drive, el observable de orden, la subarmonía primitiva de periodo
\(108T_{\rm d}\) y la estabilidad espectral finita. La realización
macroscópica material constituye un reconocimiento externo distinto y no una
laguna del resultado finito.

\section{Condiciones dinámicas de realización de un cristal temporal: Hamiltoniano, respuesta subarmónica y estabilidad}
\label{sec:programa-cristal-tiempo}

Una realización impulsada debe proporcionar, como mínimo:
\begin{enumerate}
 \item un espacio de estados físico y un Hamiltoniano autoadjunto
 \(H(t)\), o un protocolo impulsado equivalente, con
 \(H(t+T_{\rm d})=H(t)\);
 \item el propagador de Floquet
 \[
  U_F=\mathcal T\exp\!\left(
  -\frac{\mathrm i}{\hbar_{\rm int}}
  \int_0^{T_{\rm d}}H(t)\,\dd t\right);
 \]
 \item un observable de orden \(\mathcal O\) y una familia de estados para
 los que exista una respuesta subarmónica primitiva de orden \(m>1\):
 \[
  \langle\mathcal O((n+m)T_{\rm d})\rangle
  =\langle\mathcal O(nT_{\rm d})\rangle
 \]
 en el régimen asintótico, sin la misma periodicidad para un divisor propio
 de \(m\);
 \item estabilidad de \(m\), de la fase y de la amplitud frente a una clase
 abierta de perturbaciones admisibles;
 \item persistencia que no sea un transitorio de tamaño finito, junto con un
 mecanismo declarado de localización, pretermalización o protección
 dinámica.
\end{enumerate}
El calendario \(27\to54\to108\) proporciona la jerarquía interna de periodos.
La realización finita de \cref{sec:hmtf2-cristal-temporal} selecciona el orden
subarmónico \(108\) en su dominio declarado, sin sustituir las cinco
construcciones anteriores ni imponer por sí sola una fase material
macroscópica.

\begin{criterion}[Cristal de tiempo HMT--MD]
\label{crit:cristal-tiempo-hmt-md}
La identificación queda refutada si no existe un Hamiltoniano o protocolo
impulsado bien tipado; si la supuesta subarmonía desaparece bajo
perturbaciones arbitrariamente pequeñas compatibles con el modelo; si el
periodo observado es sólo el periodo impuesto por el impulso; si la
oscilación es un transitorio finito; o si el lector temporal asigna historias
equivalentes a fases físicas incompatibles.
\end{criterion}

\paragraph{Dominio matemático y físico del lector temporal.}
Las tres álgebras y sus exponenciales se demuestran en el dominio algebraico;
el lector temporal se deduce de la firma, el orden y la memoria iniciales. El
cristal temporal finito HMT satisface el
\cref{crit:cristal-tiempo-hmt-md} mediante la construcción de
\cref{sec:hmtf2-cristal-temporal}. La extensión a una fase material
macroscópica exige, además, especificar el soporte físico y pertenece al
contraste experimental posterior.

\chapter{Extensiones nulas y clausura material}
\label{chap:realizaciones-nula-material}
\index{fotón!interfaz física}\index{electrón!clausura material}
\index{extensión nula}\index{cono causal partido}

\section{\(4\) realizaciones del estado material común: nula, masiva, compuesta y topológica}
\label{sec:cuatro-modos-realizacion}

Una cifra aislada determina una coordenada, pero no identifica la estructura
que la produce, la memoria que conserva ni la operación que la convierte en
magnitud. HMT se sitúa antes de esa proyección: su objeto primario es una
extensión compatible que conserva las leyes aritméticas, la orientación
TRIT, los cocientes, la fase y la memoria de frontera. Mecánica Dimensional
especifica después qué componente se realiza como magnitud, cuál permanece
como memoria y qué transformación produce el observable.

Los desarrollos de esta parte distinguen cuatro modos por su operación de
cierre:
\begin{enumerate}
 \item una extensión abierta y nula transporta una perturbación entre
 fronteras;
 \item una extensión cerrada y persistente adquiere norma interior y define
 el candidato material;
 \item una familia de extensiones ponderada por acción define un ensamble
 térmico;
 \item una conexión que compara secciones en puntos distintos proporciona el
 codominio de la realización gravitatoria.
\end{enumerate}
La fuente común no identifica los cuatro codominios. Los mapas físicos
correspondientes se mantienen discontinuos:
\[
 \mathfrak R_\gamma:
 \Omega_{\HMT}^{\rm surv}\dashrightarrow\mathcal H_{\rm rad},
 \qquad
 \mathfrak R_m:
 \Omega_{\HMT}^{\rm surv}\dashrightarrow\mathcal H_{\rm mat},
\]
\[
 \mathfrak R_\Theta:
 \mathcal P(\Omega_{\HMT}^{\rm surv})\dashrightarrow\mathcal S_{\rm th},
 \qquad
 \mathfrak R_{\rm grav}:
 \Omega_{\HMT}^{\rm surv}\dashrightarrow\mathsf{Cartan}_{1,3}.
\]
Su unidad reside en la genealogía; su diferencia, en la lectura y en las
estructuras adicionales exigidas por cada realización.

\section{Propagación nula entre fronteras}
\label{sec:propagacion-nula-fotonica}

Sea \(B\) una base de transporte. Una \emph{sección nula} de signo
\(\pm\) es una aplicación
\[
 s_\pm:B\longrightarrow\mathcal A_{\rm sp}P_\pm
\]
tal que \(s_\pm(x)\in\mathcal A_{\rm sp}P_\pm\) para todo \(x\in B\).
Por la \cref{prop:descomposicion-espectral-partida},
\[
 N_{\rm sp}(s_\pm(x))=0
 \qquad(x\in B).
\]
Si \(\gamma\) es una ruta orientada entre dos fronteras,
\[
 \partial\gamma=x_{\rm abs}-x_{\rm em},
\]
la interfaz fotónica propuesta asigna a \(\gamma\) una sección nula
extendida. El adjetivo \emph{extendida} indica que su identidad pertenece al
transporte completo y no a un punto aislado.

La ruta conjugada se obtiene por inversión,
\[
 \gamma^\vee=C_{\rm rev}\gamma,
 \qquad
 \partial\gamma^\vee=x_{\rm em}-x_{\rm abs}.
\]
Relativamente al lector de orientación declarado, el par
\((\gamma,\gamma^\vee)\) representa las lecturas avanzada y retardada. Las
dos piernas se ordenan secuencialmente o se alojan en fibras distintas. No se
suman automáticamente dentro de una misma fibra: si se activan
simultáneamente las dos componentes conjugadas, la identidad exacta
\[
 \boxed{N_{\rm sp}(uP_++vP_-)=uv}
\]
muestra que, para \(u,v>0\), la sección abandona la frontera nula y entra en
el interior positivo.

\subsection{Longitud, ida y retorno}

Fijada la carta dimensional de velocidad,
\(\widehat c=e^{-\mathcal E}\), y una base positiva \(u_L\) de la recta de
longitud, se define
\[
 L_{\rm int}(\gamma)=\widehat c\,|\gamma|\,u_L.
\]
La inversión conserva el número de pasos:
\[
 L_{\rm int}(\gamma^\vee)=L_{\rm int}(\gamma),
 \qquad
 L_{\leftrightarrow}=2L_{\rm int}(\gamma).
\]
Estas identidades son cinemáticas. Convertirlas en propagación
electromagnética exige que un mismo mapa reproduzca la ley de Gauss, la
dinámica de calibre, la dispersión nula y la reducción a dos grados
transversales.

\subsection{Entropía de entrelazamiento entre las \(2\) lecturas}
\label{sec:entrelazamiento-ida-retorno}

La inversión de una ruta no basta para afirmar entrelazamiento cuántico. Una
formulación mínima exige primero una completación de Hilbert. Sean
\(\mathcal H_{\rm adv}\) y \(\mathcal H_{\rm ret}\) dos espacios complejos
con familias ortonormales
\(\{|i\rangle_{\rm adv}\}_{i\in I}\) y
\(\{|i\rangle_{\rm ret}\}_{i\in I}\), donde \(I\) es finito. Para una
probabilidad \(p=(p_i)_{i\in I}\), definimos
\[
 |\Psi_p\rangle
 :=\sum_{i\in I}\sqrt{p_i}\,
 |i\rangle_{\rm adv}\otimes|i\rangle_{\rm ret}.
 \tag{E1}\label{eq:estado-entrelazado-ida-retorno}
\]

\begin{proposition}[Entropía bicapa de ida y retorno]
\label{prop:entropia-entrelazamiento-ida-retorno}
Las matrices reducidas del estado
\(\rho_p=|\Psi_p\rangle\langle\Psi_p|\) son
\[
 \rho_{\rm adv}
 =\sum_{i\in I}p_i|i\rangle_{\rm adv}\langle i|,
 \qquad
 \rho_{\rm ret}
 =\sum_{i\in I}p_i|i\rangle_{\rm ret}\langle i|,
 \tag{E2}\label{eq:reducidas-ida-retorno}
\]
y poseen la misma entropía de von Neumann:
\[
 \boxed{
 S_{\rm ent}(p)
 =-\operatorname{Tr}(\rho_{\rm adv}\log\rho_{\rm adv})
 =-\sum_{i\in I}p_i\log p_i.}
 \tag{E3}\label{eq:entropia-ida-retorno}
\]
Aquí se adopta la convención \(0\log0:=0\). El estado es separable
exactamente cuando una sola probabilidad vale uno.
\end{proposition}

\begin{proof}
Al efectuar la traza parcial, la ortogonalidad anula los términos con
\(i\ne j\). Los valores propios no nulos de cada matriz reducida son, por
tanto, los \(p_i\), lo que da \eqref{eq:entropia-ida-retorno}. El rango de
Schmidt es uno exactamente en el caso determinista.
\end{proof}

La proposición formaliza, una vez elegida la completación, una entropía entre
las lecturas avanzada y retardada. No deduce las probabilidades \(p_i\), no
identifica esas lecturas con los propagadores avanzado y retardado de una
teoría de campos y no convierte una correlación clásica de rutas en
entrelazamiento. La construcción de un isomorfismo que envíe la memoria HMT
a estados físicos bipartitos, preserve la evolución y fije \(p\) pertenece a
la \textsc{interfaz física}. La entropía nula de una distribución
determinista constituye el control negativo inmediato.

\begin{statusbox}
La proposición está
\textsc{demostrada con estructura de partida explícita}: presupone la
completación de Hilbert, las bases ortonormales y la distribución \(p\). La construcción
del entrelazador que realice físicamente ambas capas es una
\textsc{interfaz física}; no forma parte de la proposición.
\end{statusbox}

\begin{definition}[Interfaz fotónica de HMT--MD]
\label{def:interfaz-fotonica-hmt-md}
Una realización fotónica de HMT--MD es una sección abierta sostenida sobre
una dirección nula del álgebra partida, transportada entre fronteras
conjugadas y protegida por una simetría que conserva exactamente la norma
nula.
\end{definition}

Los idempotentes \(P_\pm\) no son, por sí solos, fotones físicos ni las dos
polarizaciones de Maxwell. La realización completa requiere un
entrelazador
\[
 \mathfrak R_\gamma:
 \partial_{\rm null}\mathcal C_{\rm sp}^{+}
 \dashrightarrow\mathcal H_{\rm Maxwell}^{\rm phys}
\]
que preserve el producto interno, la evolución y la acción de calibre, y que
seleccione exactamente dos polarizaciones transversales en \(3+1\)
dimensiones.

\begin{criterion}[Sector fotónico]
\label{crit:sector-fotonico}
La identificación física queda refutada si la evolución abandona
\(\partial_{\rm null}\mathcal C_{\rm sp}^{+}\), genera masa sin el mecanismo
de ruptura declarado, no reproduce la dispersión nula o no deja exactamente
dos polarizaciones físicas tras la reducción de calibre.
\end{criterion}

\section{Clausura material y centro electrónico}
\label{sec:clausura-material-electron}

El acoplamiento simultáneo de dos direcciones nulas conjugadas produce, para
\(u,v>0\),
\[
 N_{\rm sp}(uP_++vP_-)=uv>0.
\]
Ninguna componente aislada posee norma distinta de cero; la norma interior
surge de su acoplamiento. Esta afirmación algebraica no suma masas
preexistentes y no identifica todavía las direcciones nulas con fotones
\emph{on shell}.

   La clasificación estructural de HMT contiene un único punto fijo de la inversión celular,
\[
 (5,5).
\]
El origen relativo de acción del electrón es \(v_e=0\). En la realización
partida normalizada se fija
\[
 Z_e=P_++P_-=I,
 \qquad
 N_{\rm sp}(Z_e)=1.
\]
El punto \((5,5)\), el origen \(v_e=0\), la firma cruda
\((24,0,12,0,-12)\) y la palabra transversal \(555555\) son cuatro objetos
distintos.

\begin{definition}[Candidato de clausura electrónica]
\label{def:candidato-clausura-electronica}
La clausura resonante de modos nulos electromagnéticos conjugados se
formaliza como el programa que busca una sección material central normalizada
obtenida por el acoplamiento de dos modos nulos conjugados, estabilizada por
una clausura de ruta y por su holonomía.
\end{definition}

La existencia de la extensión, del retorno y del centro del atlas no prueba
la compactitud mínima ni la estabilidad del cierre. El teorema dinámico
pendiente debe construir
\[
 \gamma_e=\xi_1\circ\cdots\circ\xi_N,
 \qquad
 \partial\gamma_e=0,
 \qquad
 \varepsilon(\gamma_e)=-1,
\]
donde cada \(\xi_j\) sea localmente nulo, el borde total se anule y
\(\varepsilon\) sea la holonomía de signo. La masa se asociaría al cociclo
global de clausura, no a una suma de masas nulas locales. Este programa no
reabre el electrón central \((5,5)\), el atlas ni la cadena
extensión--ruta--registro--carácter ya construidos.

\begin{proposition}[Criterio relativista de interior causal]
\label{prop:cierre-nulos-interior-causal}
Sean \(p_1,\ldots,p_N\) momentos nulos futuros en un espacio de Minkowski con
firma temporal positiva, \(p_i^2=0\), y sea
\[
 P=\sum_{i=1}^{N}p_i.
\]
Si todos los \(p_i\) son colineales, entonces \(P^2=0\). Si al menos dos son
futuros y no colineales, entonces \(P^2>0\), y puede definirse
\[
 m^2c^2=P^2.
\]
\end{proposition}

\begin{proof}
Para vectores nulos futuros se tiene
\(p_i\cdot p_j\geq0\), con igualdad exactamente cuando son colineales.
Como
\[
 P^2=\sum_i p_i^2+2\sum_{i<j}p_i\cdot p_j,
\]
la primera suma es nula. Todos los productos cruzados se anulan en el caso
colineal; si existe un par no colineal, al menos uno es estrictamente
positivo.
\end{proof}

Por tanto, una ruta cerrada no es masiva por mera topología, ni una ruta
abierta es nula por el solo hecho de poseer borde. La realización exige que
el cierre geométrico y la suma de momentos entren conjuntamente en el
interior causal.

\subsection{Möbius, conjugación y ocupación posterior}

Bajo las hipótesis de levantamiento y holonomía del
\cref{thm:dos-cuatro-pi-estadistica}, una vuelta invierte la orientación y
dos la restauran. La banda de Möbius organiza partícula y antipartícula como
orientaciones conjugadas de un mismo tipo de cierre. Su identificación física
requiere una representación que invierta la carga, preserve la norma y
entrelace la evolución. La geometría de Möbius no implica por sí sola el
álgebra de ocupación fermiónica; sus relaciones operatorias se
introducen como hipótesis operatorias en el capítulo siguiente.

\begin{criterion}[Cierre electrónico]
\label{crit:cierre-electronico}
La realización queda refutada si el centro \((5,5)\) no se eleva a una
extensión cerrada, si la holonomía no selecciona el sector impar, si el
Hamiltoniano carece de una solución ligada estable y única salvo simetría, o
si carga, dispersión y factor de forma no concuerdan con la realización
electrónica.
\end{criterion}

\chapter{Localidad de Bell y proyección holográfica}
\label{chap:localidad-bell-proyeccion}
\index{Bell, desigualdad de}\index{CHSH}\index{EPR}
\index{localidad}\index{no señalización}\index{independencia de ajustes}

La formulación EPR, el teorema de Bell y la condición de no señalización
responden a preguntas distintas \cite{EinsteinPodolskyRosen1935,Bell1964,
ClauserHorneShimonyHolt1969}. Una memoria global compartida no constituye por
sí sola una violación de localidad de Bell; tampoco una base ternaria o la
elección de \(\log 3\) como unidad de información modifican una desigualdad
de correlación. Esta sección fija los tipos necesarios y determina con
exactitud qué está probado y qué permanece abierto en HMT--MD.

\section{Escenario bipartito y premisas de localidad}
\label{sec:escenario-bell-hmt}

Considérense ajustes \(a,b\in\{0,1\}\), resultados
\(A,B\in\{-1,+1\}\) y una variable compartida
\(\lambda\in\Lambda\). Sea \(\rho\) una medida de probabilidad sobre
\(\Lambda\). La expresión
\[
 P(A,B\mid a,b)
 =\int_\Lambda
 P_A(A\mid a,\lambda)\,
 P_B(B\mid b,\lambda)\,d\rho(\lambda)
 \tag{B1}\label{eq:factorizacion-bell-local}
\]
reúne dos premisas diferentes:

\begin{enumerate}
\item \emph{factorización local}: condicionado por \(\lambda\), el resultado
de cada ala depende sólo de su ajuste local;
\item \emph{independencia de ajustes}: la medida satisface
\(d\rho(\lambda\mid a,b)=d\rho(\lambda)\).
\end{enumerate}

Se exige, además, normalización de las probabilidades y ausencia de una
selección de datos dependiente de \(a,b,A,B\). La terminología
\emph{localidad de Bell} se reservará aquí para este conjunto explícito de
hipótesis; no designa únicamente ausencia de señal superlumínica.

Definimos
\[
 E_{ab}:=\sum_{A,B=\pm1}AB\,P(A,B\mid a,b)
\]
y el funcional CHSH
\[
 S_{\rm CHSH}:=E_{00}+E_{01}+E_{10}-E_{11}.
 \tag{B2}\label{eq:funcional-chsh}
\]

\begin{theorem}[Cota CHSH bajo factorización local]
\label{thm:cota-chsh-local}
Si se satisfacen \eqref{eq:factorizacion-bell-local}, la independencia de
ajustes, la normalización y la ausencia de postselección dependiente de los
datos, entonces
\[
 \bigl|S_{\rm CHSH}\bigr|\leq2.
 \tag{B3}\label{eq:cota-chsh-dos}
\]
\end{theorem}

\begin{proof}
Para cada \(\lambda\), sean
\[
 x_a(\lambda):=\sum_{A=\pm1}A\,P_A(A\mid a,\lambda),
 \qquad
 y_b(\lambda):=\sum_{B=\pm1}B\,P_B(B\mid b,\lambda).
\]
Se tiene \(|x_a|,|y_b|\leq1\). La contribución condicionada al funcional es
\[
 s(\lambda)
 =x_0(y_0+y_1)+x_1(y_0-y_1).
\]
Por la desigualdad triangular,
\[
 |s(\lambda)|
 \leq |y_0+y_1|+|y_0-y_1|\leq2,
\]
pues \(y_0,y_1\in[-1,1]\). Integrar respecto de la misma medida \(\rho\),
independiente de \(a,b\), conserva la cota.
\end{proof}

\begin{corollary}[No señalización y su recíproco falso]
\label{cor:no-senalizacion-bell}
Las hipótesis del teorema implican
\[
 \sum_B P(A,B\mid a,b)=P(A\mid a),
 \qquad
 \sum_A P(A,B\mid a,b)=P(B\mid b).
 \tag{B4}\label{eq:no-senalizacion-bell}
\]
La no señalización, considerada aisladamente, no implica la factorización
local de \eqref{eq:factorizacion-bell-local}.
\end{corollary}

\begin{proof}
La primera identidad se obtiene sumando el factor \(P_B\), cuya
normalización vale uno; la segunda es simétrica. La implicación recíproca no
forma parte del argumento y existen distribuciones no señalizantes fuera del
politopo local.
\end{proof}

\begin{statusbox}
El teorema y el corolario son resultados
\textsc{exactos internos} de teoría de probabilidades finita. Su aplicación
a un aparato HMT--MD está
\textsc{demostrada con estructura de partida explícita} cuando se construyen
la medida, los ajustes y los mapas de resultado.
\end{statusbox}

\section{Proyección Bell del estado enriquecido común y preservación del estatuto físico}
\label{sec:estado-enriquecido-bell}

Sea \(X_{\rm enr}\) un espacio de estados enriquecidos y sean
\[
 r_A:\{0,1\}\times X_{\rm enr}\longrightarrow[-1,1],
 \qquad
 r_B:\{0,1\}\times X_{\rm enr}\longrightarrow[-1,1]
\]
dos lectores locales. Un estado puede contener memoria de frontera,
orientación, ruta y registro; toda esa información puede incorporarse a
\(\lambda=x\in X_{\rm enr}\).

\begin{theorem}[Una proyección local de un estado global sigue siendo
Bell-local]
\label{thm:no-promocion-bell-por-proyeccion}
Supóngase que:

\begin{enumerate}
\item \(x\in X_{\rm enr}\) se distribuye según una medida \(\mu\)
independiente de \(a,b\);
\item el resultado de Alice depende sólo de \((a,x)\) y el de Bob sólo de
\((b,x)\);
\item el promedio usa todas las realizaciones sin selección contextual.
\end{enumerate}

Entonces las correlaciones inducidas satisfacen
\[
 E_{ab}=\int_{X_{\rm enr}}r_A(a,x)r_B(b,x)\,d\mu(x),
\qquad
 |S_{\rm CHSH}|\leq2.
 \tag{B5}\label{eq:no-go-proyeccion-local}
\]
\end{theorem}

\begin{proof}
Se aplica el argumento del \cref{thm:cota-chsh-local} con
\(\lambda=x\), \(x_a=r_A(a,x)\) e \(y_b=r_B(b,x)\). La riqueza interna de
\(x\) no modifica la cota mientras se conserven las tres hipótesis.
\end{proof}

Este resultado impide una inferencia improcedente: una extensión global con
memoria puede explicar por qué dos lecturas comparten estructura, pero no
produce por sí sola una violación de Bell. Para obtener
\(|S_{\rm CHSH}|>2\), una realización debe abandonar o reformular al menos una
de las premisas: factorización, independencia de ajustes o muestreo no
contextual. Si la proyección depende conjuntamente de \(a\) y \(b\), deja de
ser local en el sentido de \eqref{eq:factorizacion-bell-local}; todavía habría
que probar separadamente \eqref{eq:no-senalizacion-bell}.

\section{Distinción entre localidad, funcional ternario y dependencia contextual}
\label{sec:cotejo-bell-historico}

Intervienen tres construcciones matemáticas que no deben identificarse:

\begin{enumerate}
\item la frase arquitectónica según la cual una extensión conserva memoria
global y sus lecturas locales publican proyecciones parciales;
\item el funcional ternario cíclico \(\mathcal B_3\), cuyos resultados
admiten también el valor \(0\);
\item el funcional contextual de información trítica \(I_3^{\rm ctx}\),
para el cual puede ocurrir \(I_3^{\rm ctx}(\lambda;a,b)>0\).
\end{enumerate}

El primer estrato es una interpretación holográfica compatible con el estado
enriquecido. No especifica aún una distribución bipartita
\(P_{\rm HMT}(A,B\mid a,b)\). El segundo no es el funcional CHSH
\eqref{eq:funcional-chsh}: cambia el alfabeto de resultados, el muestreo y la
normalización. La pretendida anulación por ortogonalidad no es una identidad:
de que dos ejes sean ortogonales no se sigue que una de las dos proyecciones
de un vector se anule. En efecto, para
\(v=(1,1)\), \(u(0)=(1,0)\) y
\(u(\pi/2)=(0,1)\), ambas proyecciones valen \(1\). Por consiguiente, esa
anulación no puede utilizarse para acotar \(\mathcal B_3\).

La tercera construcción relaja expresamente la independencia de ajustes. Si
\(\kappa\leq pH\), con \(H>0\), se deduce
\(p\geq\kappa/H\), no \(p\leq\kappa/H\); la desigualdad no proporciona la
cota superior de \(p\) que requeriría la conclusión propuesta para CHSH.
Asimismo, una partición binaria admisible requiere un censo exhaustivo de
políticas; sin ese censo no se obtiene la cota. Por estas dos razones, la
desigualdad informacional \(S\leq3\) no constituye una cota de Bell local bajo
las premisas estándar.

La base del logaritmo sólo cambia la unidad:
\[
 I_3(X;Y)=\frac{I_e(X;Y)}{\log 3}.
 \tag{B6}\label{eq:cambio-base-informacion-bell}
\]
En particular, ni un trit ni \(\log 3\) seleccionan una distribución, prueban
factorización o alteran la cota \(|S_{\rm CHSH}|\leq2\).
El factor \(\log3\) de la escala térmica por decisión triádica de
\cref{def:escala-termica-triadica} es una formalización distinta: cuantifica
la información de tres alternativas equiprobables y no constituye una prueba
de Bell.

Ningún rótulo algorítmico sustituye un mapa con dominio, codominio y regla
capaz de producir correlaciones CHSH. La fórmula interpretativa «ER=EPR»
requiere un entrelazador explícito entre corte geométrico, estado bipartito,
álgebra de observables y probabilidades.

La construcción global--local demuestra la existencia del estado enriquecido
y sus proyecciones. Una realización de correlaciones de Bell debe añadir un
estado bipartito, dos familias de observables locales y una medida conjunta
que reproduzca las probabilidades. Sin esos datos no se deducen una violación
CHSH, la saturación de Tsirelson ni una equivalencia ER--EPR.

\section{Criterio de realización y protocolo de contraste}
\label{sec:programa-bell-hmt}

Una promoción científica exige construir, en este orden,
\[
\begin{aligned}
X_{\rm enr}
&\longrightarrow
\bigl(\Omega_A,\Omega_B,\mu\bigr)\\
&\longrightarrow
\bigl(r_A,r_B,\mathcal A,\mathcal B\bigr)\\
&\longrightarrow
P_{\rm HMT}(A,B\mid a,b)
\longrightarrow
\bigl(E_{00},E_{01},E_{10},E_{11},S_{\rm CHSH}\bigr),
\end{aligned}
\tag{B7}\label{eq:programa-bell-hmt}
\]
donde \(\mathcal A,\mathcal B\) son las reglas de elección de ajustes. Antes
de abrir resultados experimentales deben congelarse:

\begin{enumerate}
\item la medida \(\mu\) y el dominio de estados;
\item la regla que asigna ajustes y la hipótesis, o renuncia explícita, de
independencia;
\item los mapas de resultado y el tratamiento de salidas nulas;
\item la regla de muestreo, pérdidas y coincidencias;
\item las cuatro correlaciones y la convención de signos;
\item la prueba de no señalización;
\item el intervalo de confianza y el falsador previamente registrado.
\end{enumerate}

Sólo después puede evaluarse si el mecanismo conserva la localidad de la
dinámica interna y, a la vez, produce una distribución proyectada no
Bell-local. Esta doble afirmación no se deduce del lenguaje holográfico: debe
probarse sobre \eqref{eq:programa-bell-hmt}.

\begin{criterion}[Control de Bell, localidad y no localidad]
\label{crit:bell-localidad-no-localidad}
La realización queda refutada o indebidamente tipada si ocurre cualquiera de
los siguientes hechos:

\begin{enumerate}
\item se afirma una desigualdad sin definir \(P(A,B\mid a,b)\), los ajustes y
los resultados;
\item se conserva \eqref{eq:factorizacion-bell-local} con independencia de
ajustes y, no obstante, se certifica \(|S_{\rm CHSH}|>2\);
\item se usa \(I_3(\lambda;a,b)>0\) y se denomina al modelo
«independiente de ajustes»;
\item una postselección dependiente de ajustes o resultados se oculta en la
proyección;
\item la distribución proyectada permite señalización;
\item se identifica \(\mathcal B_3\), \(\log 3\) o «ER=EPR» con CHSH sin
construir el entrelazador correspondiente.
\end{enumerate}
\end{criterion}

El censo ejecutable asociado enumera las dieciséis estrategias deterministas
locales binarias y verifica que todas satisfacen
\(|S_{\rm CHSH}|=2\). Por convexidad, este control cubre el politopo local
binario completo; no constituye una simulación de una interfaz HMT aún
ausente.

\HMTRestoreHeadings

\chapter{Información reversible, entrelazamiento y escala térmica}
\label{chap:informacion-temperatura-rev7}
\HMTClaim{MD-R-INFORMATION}
\HMTClaim{MD-R-LANDAUER-REVERSIBLE}
\HMTClaim{MD-R-BOLTZMANN}
\HMTDemoteHeadings
\input{sections/md/04g_informacion_reversible_landauer_rev7}
\input{sections/md/04e_termodinamica_rutas}
\HMTRestoreHeadings

\chapter{Funcional angular de incidencia hexada--octada y especialización de Holst}
\label{chap:barbero-rev7}
\HMTClaim{MD-R-HEXAD-OCTAD}
\HMTClaim{MD-R-BARBERO}
\HMTDemoteHeadings
\input{sections/md/07_hexada_octada_barbero}
\HMTRestoreHeadings

\chapter{Bisagra bicapa y lectores físicos hermanos}
\label{chap:bisagra-bicapa-rev8}
La incidencia hexada--octada y los canales \(q_\pm\) fijan primero el
funcional orientado \(\Gamma_{6\to8}\). Sobre la misma bicapa se lee después
la información entrelazada: \(S_{\rm bi}\) es par bajo
\(C^*\mapsto-C^*\), mientras \(\Gamma_{6\to8}\) es impar. Una lectura olvida
la orientación de las hojas y la otra la conserva; ninguna redefine
retroactivamente \(A,C^*\) o los canales.
\HMTDemoteHeadings
\input{sections/md/07a_paridad_bicapa_rev6}
\HMTRestoreHeadings

\chapter{Vacancias, residuo angular y momento magnético leptónico}
\label{chap:vacancias-radion-g2-rev7}
\HMTClaim{MD-R-RADION-RESIDUAL}
\HMTClaim{MD-R-MUON-G2}
\HMTDemoteHeadings
\input{sections/md/09_vacancias_radion_g2}
\HMTRestoreHeadings

\chapter{Transformaciones de sabor y simetrías CP y CPT}
\label{chap:ckm-cpt-rev7}
\HMTClaim{MD-R-CKM}
\HMTClaim{MD-R-CP}
\HMTClaim{MD-R-CPT}
\HMTDemoteHeadings
\chapter{Matriz de Cabibbo--Kobayashi--Maskawa y violación de carga--paridad}
\label{chap:ckm}

La matriz de Cabibbo--Kobayashi--Maskawa (CKM) entrelaza la base de sabor de
los quarks con la base de autovectores de masa en tres generaciones. La
simetría de conjugación de carga y paridad se denota por CP\@. La
transformación racional que fija los tres ángulos no se introduce como una
tabla: se construye desde el sexteto de incidencia \(\pi\)--Witt y sus reglas
sectoriales. Después se evalúa sobre las coordenadas angulares
\[
 h_{\rm CKM}=\frac{\Cstar}{6},
 \qquad
 \Delta=\frac{180}{\pi}\Dfour.
\]
Las consecuencias de invertibilidad, unitariedad y no anulación del
invariante de Jarlskog se derivan entonces de esa construcción y de las
coordenadas internas.

\section{Genealogía interna de los coeficientes de cruce}
\label{sec:genealogia-n42-ckm-rev7}

El sexteto canónico de la incidencia de nivel \(N42\) es
\[
 (b,p,q,h_6,o_8,t)=(4,5,7,6,8,3).
\]
Sus tres reglas sectoriales son los vectores fila
\[
 R_{12}=(2,-p,bq),\qquad
 R_{23}=\left(0,q,-\frac{p^2}{2}\right),
\]
\[
 R_{13}=\left(1,-pb,-\frac{o_8-h_6}{t}\right).
\]

\begin{theorem}[Construcción sectorial de la transformación CKM]
\label{thm:construccion-n42-ckm-rev7}
Dadas las reglas anteriores y el sexteto \(N42\), la única transformación
que las tiene por filas es
\[
 \boxed{
 M_{\rm CKM}=
 \begin{pmatrix}
 2&-5&28\\
 0&7&-25/2\\
 1&-20&-2/3
 \end{pmatrix}.}
\]
Los coeficientes no dependen de un ajuste angular posterior.
\end{theorem}

\begin{proof}
La sustitución directa da
\[
 bq=4\cdot7=28,\qquad
 \frac{p^2}{2}=\frac{5^2}{2}=\frac{25}{2},
\]
\[
 pb=5\cdot4=20,qquad
 \frac{o_8-h_6}{t}=\frac{8-6}{3}=\frac23.
\]
Cada regla fija sus tres entradas; por tanto las tres filas y la matriz
completa quedan determinadas.
\end{proof}

La unicidad demostrada es relativa al sistema sectorial \(N42\) y a sus
seis invariantes. La realización física exige además entrelazar el espacio
de sabor quark con el espacio de masa; ese mapa no vuelve condicional la
construcción racional anterior ni se sustituye por el contraste PDG.

\section{Transformación racional de las coordenadas angulares}

Definimos
\[
 \begin{pmatrix}\theta_{12}\\\theta_{23}\\\theta_{13}\end{pmatrix}
 =M_{\rm CKM}
 \begin{pmatrix}A\\h_{\rm CKM}\\\Delta\end{pmatrix},
 \qquad
 M_{\rm CKM}=
 \begin{pmatrix}
 2&-5&28\\
 0&7&-25/2\\
 1&-20&-2/3
 \end{pmatrix},
\]
con
\[
 \delta_{\rm CP}=9A.
\]

\begin{theorem}[Invertibilidad del modelo angular CKM]
\[
 \det M_{\rm CKM}=-\frac{3857}{6}\ne0,
\]
y
\[
 M_{\rm CKM}^{-1}=
 \begin{pmatrix}
 1528/3857&3380/3857&801/3857\\
 75/3857&176/3857&-150/3857\\
 6/551&-30/551&-12/551
 \end{pmatrix}.
\]
Por tanto, los tres ángulos y las tres coordenadas generadoras contienen la
misma información lineal.
\end{theorem}
\begin{proof}
La expansión del determinante da \(-3857/6\). La multiplicación directa de
la matriz definida por la candidata a inversa produce la identidad.
\end{proof}

Se obtiene inmediatamente la relación lineal exacta
\[
\boxed{
 3857\,\delta_{\rm CP}
 =9(1528\theta_{12}+3380\theta_{23}+801\theta_{13}).}
\]
Esta ecuación pertenece a la parametrización estándar declarada en esta
sección. El ángulo \(\delta_{\rm CP}\) depende de esa parametrización; el
invariante de Jarlskog construido más adelante es la cantidad independiente
de las redefiniciones de fase.

\section{Matriz unitaria}

Los valores angulares se expresan en grados en las relaciones lineales y en
las tablas. En toda función trigonométrica o exponencial se emplea su
representante en radianes, \(\theta_{\rm rad}=\pi\theta/180\); para aligerar
la notación se conserva el mismo símbolo para dicho argumento.

Sea \(s_{ij}=\sin\theta_{ij}\), \(c_{ij}=\cos\theta_{ij}\). En la
parametrización estándar,
\[
V=
\begin{pmatrix}
c_{12}c_{13}&s_{12}c_{13}&s_{13}e^{-i\delta}\\
-s_{12}c_{23}-c_{12}s_{23}s_{13}e^{i\delta}&
c_{12}c_{23}-s_{12}s_{23}s_{13}e^{i\delta}&s_{23}c_{13}\\
s_{12}s_{23}-c_{12}c_{23}s_{13}e^{i\delta}&
-c_{12}s_{23}-s_{12}c_{23}s_{13}e^{i\delta}&c_{23}c_{13}
\end{pmatrix}.
\]

\begin{proposition}[Unitariedad]
Para cualesquiera ángulos reales y fase real, \(V^\dagger V=I_3\).
\end{proposition}
\begin{proof}
La matriz es el producto de tres rotaciones unitarias elementales, con la fase
insertada en la rotación \(13\). El producto de matrices unitarias es
unitario.
\end{proof}

La evaluación es
\[
 (\theta_{12},\theta_{23},\theta_{13},\delta_{\rm CP})
 =(13.003313589509,\;2.398056157332,
 0.213977382244,\;65.676173123554)^\circ.
\]
Los módulos son
\[
 |V|\simeq
 \begin{pmatrix}
 0.974350259&0.225005836&0.003734601\\
 0.224873110&0.973489279&0.041841465\\
 0.008582400&0.041121745&0.999117283
 \end{pmatrix}.
\]

\section{Invariante de Jarlskog}

El invariante independiente de base es
\[
 J=c_{12}c_{23}c_{13}^2s_{12}s_{23}s_{13}\sin\delta_{\rm CP}.
\]
En el punto anterior,
\[
 \boxed{J=3.1189723326632974\times10^{-5}\ne0.}
\]

\begin{theorem}[Violación CP en la matriz construida]
Si los cuatro parámetros del modelo se insertan en el sector de tres
generaciones, \(J\ne0\); por tanto la matriz no puede hacerse real mediante
redefiniciones de fase y la matriz contiene violación CP
\cite{Cabibbo1963,KobayashiMaskawa1973,Jarlskog1985}.
\end{theorem}
\begin{proof}
Los tres senos, los tres cosenos y \(\sin\delta_{\rm CP}\) son no nulos en la
evaluación. Su producto es el valor indicado, luego el invariante de Jarlskog
no desaparece.
\end{proof}

\section{Contraste con la referencia CKM}
\label{sec:comparacion-ckm-pdg-2026}

La \cref{tab:comparacion-angular-ckm-pdg-2026} confronta la evaluación
anterior, mantenida fija, con el ajuste global de tres generaciones publicado
en la revisión CKM de 2026 del Particle Data Group
\cite{PDGCKM2026}. El ajuste externo proporciona
\(\sin\theta_{12}\), \(\sin\theta_{13}\), \(\sin\theta_{23}\) y
\(\delta\) en radianes; los ángulos en grados se obtienen mediante la
conversión monótona de la parametrización estándar. Para una incertidumbre
asimétrica, el residuo normalizado emplea la rama correspondiente al signo de
la diferencia.

\begin{table}[htbp]
\centering
\caption{Parámetros angulares CKM e invariante de Jarlskog: evaluación
HMT--MD y ajuste global PDG 2026.}
\label{tab:comparacion-angular-ckm-pdg-2026}
\small
\begin{tabular}{@{}lrrr@{}}
\toprule
Parámetro & HMT--MD & PDG 2026 & Residuo normalizado\\
\midrule
\(\theta_{12}\) (grados) &
\(13.00331359\) & \(13.01287\pm0.03999\) & \(-0.24\)\\
\(\theta_{23}\) (grados) &
\(2.39805616\) & \(2.40082^{+0.04645}_{-0.03957}\) & \(-0.07\)\\
\(\theta_{13}\) (grados) &
\(0.21397738\) & \(0.215605^{+0.005042}_{-0.004756}\) & \(-0.34\)\\
\(\delta_{\rm CP}\) (grados) &
\(65.67617312\) & \(66.1193\pm1.4324\) & \(-0.31\)\\
\(10^5J\) &
\(3.11897233\) & \(3.16^{+0.13}_{-0.11}\) & \(-0.37\)\\
\bottomrule
\end{tabular}
\end{table}

Esta tabla es \textsc{reconocimiento externo}; no interviene en la
definición de \(M_{\rm CKM}\), en la evaluación de los cuatro parámetros ni
en la prueba de unitariedad. Los cinco residuos no constituyen cinco
contrastes independientes. El ajuste global es correlacionado y está
condicionado a la unitariedad de tres generaciones; sin su matriz de
covarianza, la tabla sólo ofrece diagnósticos marginales. El contraste
multivariado científicamente vinculante sigue siendo el
\(\chi^2\) definido más adelante.

La conclusión se refiere a CP\@. La simetría discreta CPT de rutas ya
construida actúa sobre otro objeto. Su realización como teorema cuántico de
campos se compara después con las hipótesis de localidad, covariancia de
Lorentz y positividad del espectro \cite{StreaterWightman2000}.

\section{Invertibilidad y comparación estadística}

La invertibilidad de \(M_{\rm CKM}\) demuestra que alterar uno de los tres
ángulos altera de manera determinada al menos una coordenada generadora. Para
comparar el modelo con un ajuste experimental debe usarse el vector de cuatro
parámetros independientes
\[
 y=(\sin\theta_{12},\sin\theta_{23},\sin\theta_{13},\delta_{\rm CP})
\]
y su covarianza completa:
\[
 \chi^2=(y-y_0)^{\!T}\Sigma^{-1}(y-y_0).
\]
Los nueve módulos de \(V\) no son nueve observables independientes: todos son
funciones de los cuatro parámetros. Esta reducción evita contar varias veces
la misma información.

La construcción \(N42\), la inversa racional, la unitariedad y el valor no
nulo de \(J\) forman una sola cadena: primero se determinan los coeficientes,
después se evalúan los ángulos y sólo al final se abre la referencia. La
identificación del entrelazador físico sabor--masa conserva su tipo propio;
no autoriza a rebajar \(M_{\rm CKM}\) a un dato ni a obtenerla desde los
valores observados.

\input{sections/md/08a_cpt_discreto_rutas_rev6}
\HMTRestoreHeadings

\chapter{Relaciones constitutivas y transducción dimensional}
\label{chap:constitutivas-rev7}
\HMTClaim{MD-R-CONSTITUTIVES}
\HMTClaim{MD-R-DIMENSIONAL-FUNCTOR}
\HMTDemoteHeadings
\input{sections/md/delta_127/10b_funtor_dimensional}
\input{sections/md/10_constitutivas_si}
\HMTRestoreHeadings

\chapter{Holonomía, torsión y dinámica de Cartan--Holst}
\label{chap:gravedad-rev7}
\HMTClaim{MD-R-G-TYPE}
\HMTClaim{MD-R-GRAV-CARTAN}
\HMTClaim{MD-R-GRAV-PCH}
\HMTClaim{MD-R-GRAV-TORSION}
\HMTClaim{MD-R-GRAV-COLLECTIVE}
\HMTDemoteHeadings
\input{sections/md/10c_puente_retorno_holonomia_cartan}
\input{sections/md/11c_gravedad_holonomia_rev6}
\input{sections/md/11b_gravedad_sin_graviton_fundamental}
\HMTRestoreHeadings

\chapter{Memoria de frontera y reducciones cosmológicas}
\label{chap:cosmologia-rev7}
\HMTClaim{MD-R-HYDRODYNAMICS-YM}
\HMTClaim{MD-R-COSMOLOGY}
\HMTDemoteHeadings
\input{sections/md/11h_hoja_comun_hidrodinamica_yang_mills}
\input{sections/md/11i_memoria_tres_lecturas_rev6}
% Integración por delta: arquitectura cosmológica con frontera, rebote y
% componentes. No altera los teoremas Cartan--Holst ni el selector de G.

\chapter{Frontera, rebote y componentes cosmológicas}
\label{chap:cosmologia-frontera-rebote}
\index{cosmología!programa de realización}
\index{rebote cosmológico}\index{componente cosmológica descendiente}
\index{cobordismo}\index{información!conservación cosmológica}

Este capítulo distingue dos dominios. La primera sección demuestra las tres
foliaciones de de Sitter. Las secciones posteriores definen las tuplas de
datos necesarias para una realización cosmológica: selector HMT, carta de
frontera, dinámica de rebote y componente descendiente. Cada conclusión se
restringe al subdominio en el que se verifican las ecuaciones de evolución y
las condiciones de empalme declaradas.

\section{Atlas AdS--euclídeo--dS inducido por las hojas del TRIT}
\label{sec:atlas-ads-euclideo-ds-hmt}

Las tres álgebras cuadráticas del TRIT reciben una realización geométrica
coordinada. Para escalas \(\ell,H^{-1}>0\), fijamos las cartas
\[
 \mathcal M_-=(\mathbb D, g_{\rm AdS}),
 \qquad
 \mathcal M_0=(\mathscr H_0^{(L)},g_{\rm E}),
 \qquad
 \mathcal M_+=(\mathrm{dS}_4,g_{\rm dS}),
\]
donde la sección espacial hiperbólica se representa en el disco de Poincaré
por
\[
 g_{\rm AdS}^{\rm sp}
 =\frac{4\ell^2\,|dz|^2}{(1-|z|^2)^2},
 \qquad |z|<1,
\]
la L gnomónica porta la carta euclídea \(g_{\rm E}\), y la carta positiva
adopta una de las foliaciones de de Sitter construidas en la sección
siguiente. El functor geométrico conserva la etiqueta de hoja y los mapas de
cambio:
\[
 \mathfrak R_{\rm geo}:\mathcal X_{\rm TPK}^{\rm enr}
 \longrightarrow
 \mathcal M_-\cup\mathcal M_0\cup\mathcal M_+.
\]

La completación EMRH distingue interior, corona y borde, mientras la doble
proyección conserva una misma sección antes de sus evaluaciones. Sobre la
carta negativa, la restricción al borde y la representación de incidencia
forman el cuadrado
\[
\begin{CD}
 \mathcal X_{\rm TPK}^{\rm enr} @>{\mathfrak R_{\rm AdS}}>>
 \mathcal M_-\\
 @V{\mathcal P_{\rm exc}}VV @VV{\operatorname{Res}_{\partial}}V\\
 \mathscr E_{\partial} @>>{\mathfrak R_{\rm conf}}>
 \operatorname{Obs}(\partial\mathcal M_-),
\end{CD}
\]
con \(\mathfrak R_{\rm conf}\circ\mathcal P_{\rm exc}
=\operatorname{Res}_{\partial}\circ\mathfrak R_{\rm AdS}\). Esta identidad
realiza en HMT la correspondencia holográfica interior--borde: el estado del
interior y los observables conformes de frontera son dos representaciones de
la misma genealogía.

La carta euclídea centraliza el descenso local de los lectores inercial y
gravitatorio,
\[
 \mathcal G_L/\mathcal I_L=1,
\]
y constituye la residencia del principio de equivalencia. El estado anterior
a esa proyección conserva la Ambivalencia areal--volumétrica,
\[
 \mathcal L_{\rm ar}/\mathcal L_{\rm vol}
 =\pi_{\rm HMT}/\varphi_{\rm HMT}^{2}.
\]
La carta positiva completa el atlas con \(\ddot a/a=H^2>0\). Por tanto, la
misma terna geométrica contiene la correspondencia holográfica en curvatura
negativa, la equivalencia local en curvatura nula y la aceleración cosmológica
en curvatura positiva.

\section{Tres foliaciones exactas de de Sitter}
\label{sec:tres-foliaciones-de-sitter}

\begin{proposition}[Foliaciones cerrada, plana y abierta]
\label{prop:tres-foliaciones-de-sitter}
Para \(H>0\), las tres foliaciones FLRW de de Sitter de radio \(H^{-1}\)
se expresan, con \(a_\ast>0\) en la carta plana, mediante
\[
\begin{array}{ccc}
\toprule
k&a_k(t)&\text{dominio}\\
\midrule
+1&H^{-1}\cosh(Ht)&t\in\mathbb R\\
0&a_\ast e^{Ht}&t\in\mathbb R\quad\text{(parche plano)}\\
-1&H^{-1}\sinh(Ht)&t>0\quad\text{(parche abierto)}
\\\bottomrule
\end{array}
\]
y satisfacen
\[
\boxed{
\left(\frac{\dot a_k}{a_k}\right)^2+\frac{k}{a_k^2}=H^2,
\qquad
\frac{\ddot a_k}{a_k}=H^2,}
\]
\[
\boxed{
R^{(4)}=12H^2,\qquad
R^{(3)}=\frac{6k}{a_k^2}.}
\]
\end{proposition}

\begin{proof}
Para \(k=+1\), la primera ecuación utiliza
\(\tanh^2+\operatorname{sech}^2=1\). Para \(k=-1\), utiliza
\(\coth^2-\operatorname{csch}^2=1\). En la carta plana,
\(\dot a/a=H\). Las tres funciones cumplen
\(\ddot a/a=H^2\), y la fórmula FLRW de curvatura da
\(R^{(4)}=12H^2\). Cada rebanada espacial tridimensional tiene curvatura
seccional constante \(k/a_k^2\); en dimensión tres, la contracción del tensor
de curvatura da por tanto
\(R^{(3)}=3(3-1)k/a_k^2=6k/a_k^2\).
\end{proof}

El signo \(k\) tipa la curvatura de las secciones espaciales, no la
curvatura del espacio-tiempo completo: las tres cartas tienen
\(R^{(4)}>0\). La carta cerrada es global; las cartas plana y abierta son
parches.

\begin{statusbox}
    Para una ruta de retorno cerrada y la elección espacial declarada
    \(D=3\), el cociclo HMT de escala aporta los factores exactos
    \(\varphi^3,\varphi^9,\varphi^{36}\) en las profundidades
    \(9,27,108\), junto con una terna espectral de memoria y una pantalla de
    firma \((3,1)\). La selección cosmológica consiste en construir, sobre
    estado enriquecido de TPK, una medida proyectiva, un factor de escala y una ecuación de
    evolución que entrelacen esos objetos con las foliaciones anteriores. Su
    variable de duración debe derivarse de la recta de acción del
    \cref{chap:escala-accion} y de su carta metrológica, no identificarse con
    el mero número de pasos. \(H_0\), las densidades cosmológicas y el signo
    de aceleración no se emplean como entradas del selector.
\end{statusbox}

Las tres foliaciones constituyen el codominio geométrico exacto de
comparación. El programa HMT debe determinar cuál carta o atlas de cartas
realiza cada régimen, cómo se componen sus dominios y qué observables
diferencian la selección. Esta exigencia no modifica la proposición clásica
ni identifica automáticamente los tres regímenes del TRIT con
\(k=+1,0,-1\).

\section{Universo con frontera}
\label{sec:universo-con-frontera-hmt}

La completación
\[
 1000=729+271=729+243+27+1
\]
distingue interior, estratos de frontera y sello. En una realización
cosmológica, esta estratificación orienta el tipo de una construcción sobre
un par
\[
 (M,\partial M),
\]
dotado de co-tétrada, conexión de Cartan, materia, memoria y flujo. La
identidad cardinal no produce por sí sola una variedad, una métrica, una
co-tétrada ni condiciones de borde. Es una guía de organización que sólo
adquiere contenido físico mediante un mapa
\[
 \mathfrak R_{\rm cos}:
 \mathcal X_{\rm TPK}^{\rm enr}
 \dashrightarrow
 \bigl(M,\partial M,e,\omega,\Psi,\mathcal B_{\partial M}\bigr).
\]

La frontera forma parte constitutiva del problema variacional. Deben
declararse su carácter espacial, temporal o nulo; los términos de borde y de
esquina; los datos que se fijan; y las condiciones de compatibilidad con las
ecuaciones de movimiento. La memoria HMT proporciona un dominio candidato
para registrar historias y prolongaciones, pero su identificación con datos
geométricos exige el entrelazamiento de holonomías del
\cref{crit:puente-tpk-cartan}.

\section{Concentración extrema y rebote condicionado}
\label{sec:rebote-cosmologico-condicionado}

Una concentración no acotada de densidad de energía, acción o curvatura no
puede atribuirse a un mecanismo de almacenamiento energético en el sector
torsional. En la acción mínima, la corriente de espín determina localmente la
contorsión; la curvatura conserva su sector radiativo. Una corrección
repulsiva de alta densidad requiere un modelo de materia con espín, un
acoplamiento, una ecuación de estado y un signo efectivo determinados.

En un modelo de Einstein--Cartan con fluido de espín puede aparecer
\[
 s^2\propto a^{-6}.
\]
Si \(\rho_{\rm m}\) denota densidad de masa, una ecuación efectiva puede
adoptar la forma
\[
 H^2
 =
 \frac{8\pi G}{3}
 \bigl(\rho_{\rm m}-\rho_{\rm spin,m}\bigr)
 +\cdots.
\]
Si se emplea densidad de energía \(\rho_{\rm E}\), la conversión dimensional
obliga a escribir
\[
 H^2
 =
 \frac{8\pi G}{3c^2}
 \bigl(\rho_{\rm E}-\rho_{\rm spin,E}\bigr)
 +\cdots.
\]
No se mezclan ambas convenciones en una misma ecuación.

El candidato a rebote se sitúa en
\[
 \rho_{\rm m}=\rho_{\rm spin,m}
\]
o en la igualdad energética equivalente, siempre que la solución sea regular
y satisfaga las condiciones de materia y frontera declaradas. La acción
Cartan--Holst permite estudiar este mecanismo; HMT debe construir todavía el
mapa que seleccione la corriente, el umbral y el registro de información sin
consultarlos retrospectivamente.

\section{Cobordismo y componentes cosmológicas descendientes}
\label{sec:cobordismo-universo-hijo}

La apertura de una componente cosmológica descendiente se formula mediante un cobordismo
\[
 \mathcal W:
 \Sigma_{\rm in}
 \longrightarrow
 \Sigma_{\rm prog}\sqcup\Sigma_{\rm desc}.
\]
Antes de afirmar conservación de información debe fijarse el objeto
conservado. Si \(\mathscr I\) es una carga aditiva definida por un cociclo,
puede ensayarse
\[
 \mathscr I(\Sigma_{\rm in})
 =
 \mathscr I(\Sigma_{\rm prog})
 +\mathscr I(\Sigma_{\rm desc})
 +\mathscr F_{\partial\mathcal W}.
\]
Si \(\mathscr I\) representa información cuántica, la estructura pertinente
es una isometría o una evolución unitaria y el balance no se reduce, en
general, a una suma escalar. Han de calcularse las entropías de los
subsistemas, sus correlaciones y la información mutua.

Una realización debe especificar:
\begin{enumerate}
 \item las condiciones de unión entre las secciones inicial, progenitora y
 descendiente;
 \item los términos de borde y esquina de la acción;
 \item la causalidad, la hiperbolicidad y el problema de valores iniciales;
 \item el control de singularidades, curvas temporales cerradas y condiciones
 de energía;
 \item la carga, isometría o ley de conservación concreta que representa la
 información.
\end{enumerate}
Las identidades de Bianchi y el balance de Noether son condiciones locales
necesarias; no demuestran por sí solos el cambio de topología ni la
conservación global de información.

\section{Componentes paralelas y modo intercomponente}
\label{sec:componentes-cosmologicas-paralelas}

Una frontera desconectada y una colección de sectores poseen composiciones
distintas. En una realización de tipo topológico,
\[
 \mathcal H_{\Sigma_{\rm parent}\sqcup\Sigma_{\rm child}}
 \simeq
 \mathcal H_{\rm parent}\otimes\mathcal H_{\rm child}
\]
representa la composición de componentes de una frontera. En cambio,
\[
 \mathcal H_{\rm cos}
 =\bigoplus_a\mathcal H_a
\]
representa alternativas o sectores de superselección. El producto tensorial
y la suma directa no se sustituyen mutuamente.

Sobre la suma directa, un operador global candidato tiene forma de bloques:
\[
 \mathbb H_{\rm cos}
 =
 \begin{pmatrix}
 H_1&K_{12}&\cdots\\
 K_{21}&H_2&\cdots\\
 \vdots&\vdots&\ddots
 \end{pmatrix}.
\]
Para definir un operador físico deben fijarse un dominio denso común o
hipótesis de acotación relativa, un reloj común o relacional y la condición
\[
 K_{ba}=K_{ab}^{\dagger}.
\]
Los bloques no diagonales han de derivarse de la frontera, del cobordismo o
de otro acoplamiento declarado. Si son no nulos, las componentes forman
sectores acoplados de una teoría global; no son universos independientes en
sentido fuerte.

Un vector propio con soporte en varios sumandos constituye un modo
intercomponente. Denominar gravitón a su cuanto exige, además, que el sector
físico reducido posea helicidades \(\pm2\), dos polarizaciones transversas,
energía positiva, propagación lumínica, ausencia de fantasmas y taquiones,
identidades de Ward y acoplamiento universal. El soporte global por sí solo
no satisface estas condiciones.

\section{Condiciones de cierre y refutación}
\label{sec:condiciones-prueba-cosmologia-hmt}

La arquitectura se promoverá únicamente si se construyen:
\begin{enumerate}
 \item soluciones regulares de rebote con materia y signo derivados;
 \item un problema de frontera bien planteado;
 \item un cobordismo admisible con condiciones de unión;
 \item un transporte de información definido como carga, isometría o
 evolución unitaria;
 \item un operador global autoadjunto con acoplamientos no diagonales
 seleccionados antes de la comparación;
 \item predicciones que distingan la realización de una cosmología de una
 sola componente.
\end{enumerate}

La propuesta falla si no existe la solución regular, si las condiciones de
unión son incompatibles, si se viola la carga o isometría escogida, si el
cobordismo depende de una carta arbitraria, si el operador global resulta
necesariamente diagonal o si la evolución pierde causalidad o
hiperbolicidad sin un mecanismo alternativo explícito. No se invoca una
energía global canónica sin simetría temporal ni definición cuasilocal.

\input{sections/md/11e_finitud_anisotropia}
\chapter{Límite central, órbitas de Kepler y holonomía de Sommerfeld}
\label{chap:limite-central-kepler-sommerfeld}
\index{Kepler, órbita de}\index{Sommerfeld, cuantización de}
\index{Levi-Civita, cubierta de}\index{Möbius, holonomía de}

Este capítulo distingue cuatro estratos que la terminología física puede
superponer indebidamente: la cancelación exacta de flujos interiores en un
grafo; la realización continua de un campo central; la dinámica kepleriana
que resulta del potencial inverso; y la condición semiclásica de fase sobre
una banda con holonomía. La primera es combinatoria. Las tres restantes
requieren, sucesivamente, una carta métrica, un lagrangiano y una línea de
fase semiclásica.

\section{Gauss discreto y descenso inverso al área}
\label{sec:gauss-discreto-limite-central}

Sea \(G=(V,E)\) un grafo finito no orientado, sea
\(\vec E:=\{(x,y),(y,x):\{x,y\}\in E\}\) su conjunto de arcos y sea
\(F:\vec E\to\mathbb R\) una \(1\)-co\-cadena antisimétrica,
\(F(y,x)=-F(x,y)\). Definimos
\[
 \operatorname{div}F(x)
 :=\sum_{y:\{x,y\}\in E}F(x,y).
\]
Para \(S\subseteq V\), el flujo
\(\operatorname{Flux}_{\partial S}F\) es la suma orientada sobre las aristas
con un extremo en \(S\) y el otro en \(V\setminus S\).

\begin{theorem}[Teorema de Gauss en grafos]
\label{thm:gauss-grafo-hmt-md}
Para todo \(S\subseteq V\),
\[
 \sum_{x\in S}\operatorname{div}F(x)
 =\operatorname{Flux}_{\partial S}F.
 \tag{K1}\label{eq:gauss-grafo-hmt-md}
\]
\end{theorem}

\begin{proof}
Cada arista con ambos extremos en \(S\) aparece dos veces con orientaciones
opuestas y se cancela. Sólo permanecen las aristas que atraviesan la cadena
de borde.
\end{proof}

En la red cúbica, para \(R\in\mathbb Z_{\geq0}\),
\[
 C_R:=\{x\in\mathbb Z^3:\|x\|_\infty\leq R\},
\]
el número de enlaces salientes es
\[
 |\partial C_R|=6(2R+1)^2.
\]
Si una fuente posee flujo total \(Q\) y se impone uniformidad media sobre el
cascarón, el flujo por enlace vale
\[
 E_R=\frac{Q}{6(2R+1)^2}.
\]
Al definir \(r_{\rm eff}\) mediante
\[
 4\pi r_{\rm eff}^2=6(2R+1)^2,
\]
se obtiene la identidad exacta
\[
 E_R=\frac{Q}{4\pi r_{\rm eff}^2}.
 \tag{K2}\label{eq:inverso-area-cascaron}
\]
El paso desde \eqref{eq:gauss-grafo-hmt-md} a
\eqref{eq:inverso-area-cascaron} usa la hipótesis adicional de isotropía
media. Identificar \(Q\) con una fuente gravitatoria o eléctrica requiere,
además, el transductor dimensional correspondiente.

\section{Lagrangiano central y cónicas de Kepler}
\label{sec:lagrangiano-central-kepler}

Sea \(\mu>0\) una masa reducida y
\(\mathcal K>0\) un acoplamiento del tipo producto de energía y longitud,
\([\mathcal K]=\mathsf E\,\mathsf L\). En la hoja
euclídea continua consideremos
\[
 L_{\rm Kep}(\mathbf r,\dot{\mathbf r})
 =\frac{\mu}{2}\|\dot{\mathbf r}\|^2+\frac{\mathcal K}{r},
 \qquad
 H_{\rm Kep}(\mathbf r,\mathbf p)
 =\frac{\|\mathbf p\|^2}{2\mu}-\frac{\mathcal K}{r},
 \tag{K3}\label{eq:lagrangiano-hamiltoniano-kepler}
\]
donde \(r=\|\mathbf r\|>0\) y
\(\mathbf p=\mu\dot{\mathbf r}\). La ecuación de Euler--Lagrange es
\[
 \mu\ddot{\mathbf r}
 =-\frac{\mathcal K}{r^3}\mathbf r.
 \tag{K4}\label{eq:ecuacion-central-kepler}
\]
El momento angular específico
\(\mathbf h=\mathbf r\times\dot{\mathbf r}\) es constante; si
\(\mathbf h\ne0\), la trayectoria permanece en el plano ortogonal a
\(\mathbf h\).

\begin{proposition}[Cónica kepleriana]
\label{prop:conica-kepleriana-hmt-md}
Sea \(h=\|\mathbf h\|\). Toda solución no radial de
\eqref{eq:ecuacion-central-kepler} satisface
\[
 r(\phi)
 =\frac{p}{1+e\cos(\phi-\phi_0)},
 \qquad
 p=\frac{\mu h^2}{\mathcal K}.
 \tag{K5}\label{eq:conica-kepleriana}
\]
Para energía negativa, \(0\leq e<1\) y la órbita es elíptica.
\end{proposition}

\begin{proof}
Con \(u=1/r\), la ecuación de Binet adopta la forma
\[
 u''(\phi)+u(\phi)=\frac{\mathcal K}{\mu h^2}.
\]
Su solución general es
\[
 u=\frac{\mathcal K}{\mu h^2}
 \bigl(1+e\cos(\phi-\phi_0)\bigr).
\]
Invertir \(u\) da \eqref{eq:conica-kepleriana}; la condición de ligadura
equivale a \(e<1\).
\end{proof}

La velocidad areolar es
\[
 \frac{\mathrm dA}{\mathrm dt}=\frac h2.
 \tag{K6}\label{eq:segunda-ley-kepler}
\]
Para una elipse de semieje mayor \(a\), integrar
\eqref{eq:segunda-ley-kepler} durante un periodo produce
\[
 A_{\rm ell}=\pi ab,
 \qquad
 b=a\sqrt{1-e^2},
 \qquad
 p=a(1-e^2)=\frac{\mu h^2}{\mathcal K}.
\]
Por tanto,
\(T=2A_{\rm ell}/h\), y al eliminar \(b\), \(p\) y \(h\) se obtiene
\[
 T^2=\frac{4\pi^2\mu}{\mathcal K}\,a^3.
 \tag{K7}\label{eq:tercera-ley-kepler}
\]
En el problema gravitatorio de prueba,
\(\mathcal K=G M\mu\), y
\eqref{eq:tercera-ley-kepler} se reduce a
\(T^2=4\pi^2a^3/(GM)\) \cite{Kepler1609,Newton1687}.

El teorema de Gauss del grafo es exacto. Bajo isotropía y distancia
declaradas, el flujo produce el campo inverso al área; el límite continuo
transporta después ese campo al lagrangiano central y a las órbitas
keplerianas escritas arriba.

\section{Cubierta de Levi--Civita y banda de Möbius}
\label{sec:levi-civita-mobius-kepler}

En el plano complejo sin el origen, la aplicación de Levi--Civita
\[
 q=z^2
 \tag{K8}\label{eq:cubierta-levi-civita}
\]
es una doble cubierta. Una vuelta de \(q\) alrededor del origen intercambia
los dos levantamientos \(z\) y \(-z\); hacen falta dos vueltas para cerrar el
levantamiento. Es
un testigo exacto de memoria \(\mathbb Z_2\). No es, sin embargo, una banda
de Möbius: la regularización completa del problema de Kepler exige también
reparametrizar el tiempo y transformar los momentos.

Sea ahora \(\gamma:S^1\hookrightarrow\mathbb R^2\) un embedding cuya imagen
es una órbita elíptica, y fijemos \(\varepsilon>0\). Definimos
\[
 M_\gamma
 :=([0,2\pi]\times[-\varepsilon,\varepsilon])/
 \bigl((0,u)\sim(2\pi,-u)\bigr).
 \tag{K9}\label{eq:banda-mobius-orbita}
\]

\begin{theorem}[Periodo orientado de la banda orbital]
\label{thm:periodo-orientado-banda-orbital}
\(M_\gamma\) es una banda de Möbius cuyo núcleo es \(\gamma\). Una sección
de la línea de signo posee holonomía
\[
 \mathcal H(\gamma)=-I,
 \qquad
 \mathcal H(\gamma^2)=I.
 \tag{K10}\label{eq:holonomia-banda-orbital}
\]
Por tanto, su retorno como sección orientada requiere dos vueltas.
\end{theorem}

\begin{proof}
La identificación de las fibras extremas mediante \(u\mapsto-u\) es la
construcción estándar de la banda de Möbius. La holonomía de una vuelta es
la involución de la fibra; su cuadrado es la identidad.
\end{proof}

La cubierta \eqref{eq:cubierta-levi-civita} y la banda
\eqref{eq:banda-mobius-orbita} comparten una monodromía de orden dos, pero no
son objetos isomorfos. Esta distinción impide deducir una orientación de
partícula--antipartícula a partir de la regularización de Levi--Civita sin un
mapa adicional.

\section{Condición semiclásica de Sommerfeld con holonomía}
\label{sec:sommerfeld-holonomia-banda}

Sea
\[
 J_\gamma:=\oint_\gamma\mathbf p\cdot\mathrm d\mathbf q
\]
la acción de un ciclo y supóngase que el transporte semiclásico aporta la
fase \(\exp(\mathrm iJ_\gamma/\hbar_{\rm int})\). Si la línea real asociada
posee holonomía \(\varepsilon_\gamma\in\{+1,-1\}\), la condición de cierre de
la sección complejificada es
\[
 \varepsilon_\gamma
 \exp\!\left(\frac{\mathrm iJ_\gamma}{\hbar_{\rm int}}\right)=1.
 \tag{K11}\label{eq:cierre-fase-holonomia-sommerfeld}
\]

\begin{proposition}[Desplazamiento de media unidad por holonomía de signo]
\label{prop:desplazamiento-sommerfeld-holonomia}
Bajo las hipótesis anteriores,
\[
 \begin{array}{lll}
 \varepsilon_\gamma=+1
 &\Longrightarrow&
 J_\gamma=2\pi\hbar_{\rm int}n,\\[1mm]
 \varepsilon_\gamma=-1
 &\Longrightarrow&
 J_\gamma=2\pi\hbar_{\rm int}\left(n+\dfrac12\right),
 \end{array}
 \qquad n\in\mathbb Z.
 \tag{K12}\label{eq:cuantizacion-sommerfeld-holonomia}
\]
\end{proposition}

\begin{proof}
La \cref{eq:cierre-fase-holonomia-sommerfeld} exige una fase igual a uno en
el sector trivial e igual a menos uno en el sector de signo. Sus argumentos
difieren, respectivamente, de un múltiplo entero o semientero de \(2\pi\).
\end{proof}

La proposición es una consecuencia exacta de la fase y de la holonomía
declaradas; no selecciona por sí sola el Hamiltoniano de Coulomb, un número
cuántico físico ni un espectro atómico \cite{Sommerfeld1916}. En una
cuantización EBK deben añadirse los índices de Maslov de los ciclos
correspondientes, sin contar dos veces el desplazamiento topológico.

Relativamente al lector de orientación, las dos hojas locales pueden
representar las rutas \(\gamma\) y \(\gamma^\vee\) y enlazarse con las
dos orientaciones del entrelazamiento bicapa definido en
\cref{def:entropia-bicapa-rev6}. Esta correspondencia no identifica
esas rutas con funciones de Green avanzadas o retardadas ni afirma que la
entropía cause la cónica kepleriana: la órbita procede de
\eqref{eq:ecuacion-central-kepler}; la entropía mide, después de una
completación de Hilbert, la correlación entre hojas.

\begin{criterion}[Límite Kepler--Sommerfeld]
\label{crit:limite-kepler-sommerfeld}
La realización queda refutada o indebidamente tipada si se obtiene una ley
inversa al cuadrado sin declarar isotropía; si
\(\mathcal K\) o \(G\) se introducen como escalares sin recta dimensional;
si la cubierta \(q=z^2\) se identifica con una banda de Möbius; si
\eqref{eq:cuantizacion-sommerfeld-holonomia} se publica como espectro exacto
sin Hamiltoniano, ciclos e índices de Maslov; o si se atribuye
entrelazamiento a una mera inversión determinista de ruta.
\end{criterion}

\HMTRestoreHeadings

\chapter{Interfaz física de las reducciones de rango, dualidad T y teoría M}
\label{chap:interfaz-m-rev7}
\HMTClaim{MD-R-M-THEORY-INTERFACE}
Las pantallas exactas de la primera parte se prolongan aquí hacia su posible
lectura física. El capítulo distingue las reducciones matemáticas ya
demostradas del mapa adicional necesario para campos, cuerdas, CFT y
supergravedad undecadimensional.
\HMTDemoteHeadings
\input{sections/md/12c_interfaz_teoria_m_rev8}
\HMTRestoreHeadings

\chapter{Síntesis dimensional}
\label{chap:sintesis-md-rev7}
\HMTClaim{MD-R-SYNTHESIS}
\HMTDemoteHeadings
\input{sections/md/12b_sintesis_realizaciones}
\input{sections/md/12_sintesis_rev6}
\HMTRestoreHeadings
